\section{The mirror lemma and the mirror reduction (Theorem~A of \texorpdfstring{\dirn{MIRRORETV}}{MIRRORETV})}\label{mm:sec}
% Local typesetting setting, closed by the \endgroup at the end of this file: long Lean names in
% statement heads cannot be hyphenated, so allow loose lines instead of lines running into the margin.
\begingroup\setlength{\emergencystretch}{0.5\linewidth}

This section gives written proofs for the part of the Lean development \lean{MirrorOdd}
(file \dirn{MIRRORETV/lean/MirrorMain.lean}) that concerns point sets symmetric under a reflection:
the mirror lemma, the reduction of a mirror-symmetric set to an $x$-sorted configuration (``Theorem~A''
of \dirn{MIRRORETV/README.md}, in its Lean form), the transport from an arbitrary mirror line to the
$x$-axis, and the bounds for $k=4,5,6$. In this section ``Theorem~A'' always means this mirror reduction.
The main theorem of \dirn{WRITEUP/SYMMETRIC.md}, also called Theorem~A there, is a different statement,
about sets of at least $33$ points with a symmetry; this document states it only in conditional forms,
whose hypotheses are case classes (Theorems~\ref{sr:thm:main}, \ref{ar:thm:main} and~\ref{mb:thm:main3}).
\begin{itemize}
\item \emph{Mirror lemma} (Theorems~\ref{mm:thm:mirror-polygon} and~\ref{mm:thm:axis-polygon}). Let
  $c_0,\dots,c_{m-1}$ have strictly increasing abscissae, and let every triple $c_i,c_j,c_l$ with $i<j<l$
  be clockwise. If $m\ge2$ and all $c_i$ lie strictly above the $x$-axis, then the points $c_i$ and their
  mirror images $\operatorname{mir}(c_i)$ contain a convex $2m$-gon. If $m\ge3$ and the first (or the
  last) point lies on the axis while the others lie strictly above it, they contain a convex
  $(2m-1)$-gon. As in Lean, the proof lists the points in an explicit cyclic order and checks that
  \emph{every} triple of the list, taken in list order, is counterclockwise.
\item \emph{Theorem~A} (Theorems~\ref{mm:thm:odd}, \ref{mm:thm:even} and~\ref{mm:thm:line}). Let $j\ge2$,
  and let $M$ be a natural number such that every $x$-sorted configuration with no $(j+1)$-cap and no
  convex $(2j+1)$-gon has at most $M$ points. Then every finite set in general position that is
  symmetric under the reflection in a line and has no convex $(2j+1)$-gon has at most $2M$ points. In
  the even case ($j$-caps, convex $2j$-gons) the bound is $2M+2$. Here $M$ is an \emph{assumed} upper
  bound: the bound on configurations is a hypothesis of these theorems, and it is not proved in this
  section.
\item \emph{Small $k$} (Theorem~\ref{mm:thm:small}). Without further hypotheses, a finite set in general
  position that is symmetric in a line and has no convex $k$-gon has at most $4$, $8$, $12$ points for
  $k=4,5,6$.
\end{itemize}

\paragraph{Sources.}
{\raggedright
The files are the following (SHA-256, first 12 hexadecimal digits).
\begin{itemize}
\item The declarations treated here: \dirn{MIRRORETV/lean/MirrorMain.lean} (2677 lines,
  \texttt{ca623e55a1b9}\dots, namespace \lean{MirrorOdd}) and, for configurations and caps,
  \dirn{ErdosSzekeres/Omission/Realization.lean} (\texttt{ae4970e0808e}\dots).
\item The positive-cycle lemma (Lemma~\ref{fd:lem:pos-cycle}): \dirn{ErdosSzekeres/ConvexCycleCuts.lean}
  (\texttt{b4ad0ba9e583}\dots).
\item The basic definitions recalled in the foundations section: \dirn{ErdosSzekeres/Definitions.lean}
  (\texttt{2425500eefd1}\dots), \dirn{ErdosSzekeres/Basic.lean} (\texttt{13667aa2d3dc}\dots),
  \dirn{ErdosSzekeres/Orientation.lean} (\texttt{f67955ac7091}\dots) and
  \dirn{ErdosSzekeres/OrientationGeometry.lean} (\texttt{c7db9bbbb65e}\dots).
\end{itemize}
Lines 22--1506 of
\dirn{MIRRORETV/lean/MirrorMain.lean} hold five positive-cycle helpers and nine small polygons
(Lemma~\ref{mm:lem:cyc} and Proposition~\ref{mm:prop:small-polygons}). The file header (lines 4--10)
describes the polygons for caps of four points, says that every orientation inequality is closed by
\lean{linarith} with a certificate found offline by a script \dirn{gen_certs.py}, and lists
\dirn{MirrorPolygons.lean} and \dirn{small_polygons.lean}, generated by a script \dirn{gen_lean.py},
among the source files of the assembled file; the comment at line 993 marks the polygons for caps of two
and three points as generated. Neither script is included in \dirn{MIRRORETV}, and the header lists the
other source files without saying how they were written. Toolchain: Lean~4.24.0, Mathlib \texttt{f897ebcf}. The text of \dirn{MIRRORETV/lean/MirrorMain.lean},
lines 3--2677, occurs verbatim as lines 1562--4236 of \dirn{SYMLEAN/mirror/ReductionThreeClasses.lean}
(\texttt{0cc83994fe8c}\dots): a declaration at line $L$ of \dirn{MIRRORETV/lean/MirrorMain.lean} is at
line $L+1559$ there. That copy has its own record: lines 1--16 of \dirn{SYMLEAN/mirror/axioms.out} print
the same three axioms for the sixteen theorems listed in the Rigor paragraph below, and
\dirn{SYMLEAN/mirror/README.md} reports a compile without errors. With this record in place of
\dirn{MIRRORETV/lean/MirrorMain.out}, everything below applies to both copies.
\dirn{MIRRORETV/lean/FarFlat.lean} (\texttt{08a79f881eea}\dots) and the four Lean files of
\dirn{MIRRORETV/review/} are not used here (see the Coverage paragraph).\par}

\paragraph{Rigor.}
Every statement below that carries a Lean reference is \emph{Lean-checked}, except the two parts marked
as written steps (part~(e) of Lemma~\ref{mm:lem:orient-mir} and part~(b) of Lemma~\ref{mm:lem:slope});
Lemma~\ref{mm:lem:upper-caps} is a written step and carries no Lean reference. The recorded compile log
\dirn{MIRRORETV/lean/MirrorMain.out}, lines 848--863, prints for the sixteen theorems \lean{mirror_four},
\lean{mirror_five}, \lean{mirror_six}, \lean{mirror_seven}, \lean{mirror_eight}, their five
\lean{_line} versions, \lean{mirror_odd_general}, \lean{mirror_even_general},
\lean{mirror_odd_general_line}, \lean{mirror_even_general_line}, \lean{mirror_seven'} and
\lean{mirror_eight'} (all in \lean{MirrorOdd}) exactly the axioms \lean{propext},
\lean{Classical.choice} and \lean{Quot.sound}; the log reports no error, only 121 linter warnings about
unused simp arguments. Every other declaration treated in this section is used, directly or indirectly,
by one of these sixteen theorems, so it depends on no further axiom. The file contains no \lean{sorry},
\lean{admit}, \lean{native_decide} or \lean{axiom} declaration; in the declarations treated here
\lean{decide} only compares explicit elements of finite index types (kernel evaluation). This section
relies on the recorded log; Lean was not re-run for it. The proofs below are written proofs (not
refereed). Theorems~\ref{mm:thm:odd} and~\ref{mm:thm:even} and parts (b), (c) of Theorem~\ref{mm:thm:line} take a
parametric hypothesis (hM): an assumed bound $M$, not an unproved claim of the project; several instances of
(hM) are discharged by Lean-checked theorems (Remark~\ref{mm:rem:hm}).

\paragraph{Dependencies.}
We use the conventions and facts of the foundations section: Definitions~\ref{fd:def:point},
\ref{fd:def:orient}, \ref{fd:def:gp}, \ref{fd:def:cp}, \ref{fd:def:ngon} and~\ref{fd:def:inv}, and
Lemmas~\ref{fd:lem:orient}, \ref{fd:lem:affine-orient}, \ref{fd:lem:collinear}, \ref{fd:lem:cp-hull},
\ref{fd:lem:mono}, \ref{fd:lem:affine-cp} and~\ref{fd:lem:pos-cycle}. The bounds for $k=7,8,9$, which
combine the results of this section with the repository's Lean formalization of Baek's theorem (for
$k=9$ also with the vendored cups--caps theorem), are Theorems~\ref{mb:thm:seven}, \ref{mb:thm:eight}
and~\ref{mb:thm:nine}.

\paragraph{Conventions.}
Points are $p=(p_0,p_1)\in\R^2$ (Lean \lean{Point}, Definition~\ref{fd:def:point}); $p_0$ is the abscissa
and $p_1$ the ordinate. For $m\in\N$ we write $[m]=\{0,1,\dots,m-1\}$ with its usual order; it stands for
Lean's \texttt{Fin m}. All point sets are finite (Lean \texttt{Finset Point}). A finite set $X\subset\R^2$
\emph{has a convex $n$-gon} if \lean{HasConvexNGon} $n$ $X$ holds (Definition~\ref{fd:def:ngon}): some
$T\subseteq X$ with $\card T=n$ is in convex position. A finite set $T$ is \emph{$x$-injective} if
$p\mapsto p_0$ is injective on $T$ (Lean \texttt{Set.InjOn (fun p => p 0) T}). In Lean, $2m-1$, $m-1$
and $K-1$ in the statements below are truncated subtractions of natural numbers; wherever they occur,
the hypotheses ($m\ge3$, $K\ge4$) make them equal to the ordinary differences.

\subsection{Reflection, configurations and sorting}

\begin{definition}[reflection, upper part, axis part; \leanref{MirrorOdd.mir}{MIRRORETV/lean/MirrorMain.lean:16}, \leanref{MirrorOdd.Rset}{MIRRORETV/lean/MirrorMain.lean:1598}, \leanref{MirrorOdd.Aset}{MIRRORETV/lean/MirrorMain.lean:1599}]\label{mm:def:mir}
For $p\in\R^2$ let $\operatorname{mir}(p)=(p_0,-p_1)$, the reflection in the $x$-axis. For a finite set
$S\subset\R^2$ let
\[ R(S)=\{p\in S: p_1>0\},\qquad A(S)=\{p\in S: p_1=0\} \]
(Lean \lean{Rset} and \lean{Aset}). In all theorems below, symmetry of $S$ in the $x$-axis is the
hypothesis ``$\operatorname{mir}(p)\in S$ for every $p\in S$'' (``maps into'', as in
Definition~\ref{fd:def:inv}). Since $\operatorname{mir}$ is injective and $S$ is finite, this is
equivalent to $\operatorname{mir}(S)=S$, but only the stated form is used.
\end{definition}

\begin{lemma}[identities for $\operatorname{mir}$ and $\orient$; \leanref{MirrorOdd.mir_zero}{MIRRORETV/lean/MirrorMain.lean:18}, \leanref{MirrorOdd.mir_one}{MIRRORETV/lean/MirrorMain.lean:19}, \leanref{MirrorOdd.mir_mir}{MIRRORETV/lean/MirrorMain.lean:1543}, \leanref{MirrorOdd.orient_mir3}{MIRRORETV/lean/MirrorMain.lean:2187}, \leanref{MirrorOdd.orient_mir2}{MIRRORETV/lean/MirrorMain.lean:2190}, \leanref{MirrorOdd.orient_cycl}{MIRRORETV/lean/MirrorMain.lean:2193}, \leanref{MirrorOdd.orient_swap12}{MIRRORETV/lean/MirrorMain.lean:2196}, \leanref{MirrorOdd.orient_rev}{MIRRORETV/lean/MirrorMain.lean:2199}]\label{mm:lem:orient-mir}
For all $p,a,b,d\in\R^2$:
\begin{enumerate}[label=(\alph*)]
\item $\operatorname{mir}(p)_0=p_0$, $\operatorname{mir}(p)_1=-p_1$ and $\operatorname{mir}(\operatorname{mir}(p))=p$
  (\lean{mir_zero}, \lean{mir_one}, \lean{mir_mir});
\item $\orient(\operatorname{mir}a,\operatorname{mir}b,\operatorname{mir}d)=-\orient(a,b,d)$ (\lean{orient_mir3});
\item $\orient(\operatorname{mir}a,\operatorname{mir}b,d)=-\orient(a,b,\operatorname{mir}d)$ (\lean{orient_mir2});
\item $\orient(a,b,d)=\orient(b,d,a)$, $\orient(a,b,d)=-\orient(b,a,d)$ and $\orient(a,b,d)=-\orient(d,b,a)$
  (\lean{orient_cycl}, \lean{orient_swap12}, \lean{orient_rev});
\item (written step; not a separate Lean declaration)
  $\orient(a,b,\operatorname{mir}d)=\orient(a,b,d)-2d_1(b_0-a_0)$.
\end{enumerate}
\end{lemma}

\begin{proof}
(a) is the definition of $\operatorname{mir}$. For (b) and (e) expand
$\orient(u,v,w)=(v_0-u_0)(w_1-u_1)-(v_1-u_1)(w_0-u_0)$ (Definition~\ref{fd:def:orient}):
\begin{align*}
\orient(\operatorname{mir}a,\operatorname{mir}b,\operatorname{mir}d)
  &=(b_0-a_0)(-d_1+a_1)-(-b_1+a_1)(d_0-a_0)=-\orient(a,b,d),\\
\orient(a,b,\operatorname{mir}d)
  &=(b_0-a_0)(-d_1-a_1)-(b_1-a_1)(d_0-a_0)=\orient(a,b,d)-2d_1(b_0-a_0).
\end{align*}
(c) By (a), $d=\operatorname{mir}(\operatorname{mir}d)$, so (b) applied to $a,b,\operatorname{mir}d$ gives
$\orient(\operatorname{mir}a,\operatorname{mir}b,d)
=\orient(\operatorname{mir}a,\operatorname{mir}b,\operatorname{mir}(\operatorname{mir}d))
=-\orient(a,b,\operatorname{mir}d)$.
(d) These are the cases of Lemma~\ref{fd:lem:orient}(d) for the cyclic permutation
$(a,b,d)\mapsto(b,d,a)$ (sign $+1$), the transposition of the first two arguments and the transposition
of the first and the third argument (sign $-1$). Each can also be checked by expanding.
\end{proof}

\begin{definition}[configurations, caps and sorting; \leanref{ErdosSzekeres.Omission.Configuration}{ErdosSzekeres/Omission/Realization.lean:9}, \leanref{ErdosSzekeres.Omission.Configuration.HasCap}{ErdosSzekeres/Omission/Realization.lean:64}, \leanref{MirrorOdd.xset}{MIRRORETV/lean/MirrorMain.lean:1549}, \leanref{MirrorOdd.emb}{MIRRORETV/lean/MirrorMain.lean:1553}, \leanref{MirrorOdd.sortPt}{MIRRORETV/lean/MirrorMain.lean:1558}, \leanref{MirrorOdd.cfg}{MIRRORETV/lean/MirrorMain.lean:1584}]\label{mm:def:config}
The following notions concern finite sequences of points sorted by abscissa.
\begin{enumerate}[label=(\alph*)]
\item Let $n\in\N$. A \emph{configuration of $n$ points} is a map $C:[n]\to\R^2$, $i\mapsto C_i$ (Lean
  field \lean{point}), such that $i\mapsto(C_i)_0$ is strictly increasing (field \lean{xOrder}) and
  $\orient(C_i,C_j,C_k)\ne0$ for all $i<j<k$ in $[n]$ (field \lean{generalPosition}). Its \emph{point
  set} is $\mathrm{pts}(C)=\{C_i:i\in[n]\}$ (Lean \texttt{Finset.univ.image c.point}).
\item Let $k\in\N$. A configuration $C$ of $n$ points \emph{has a $k$-cap} (Lean \texttt{c.HasCap k}) if
  there is a strictly increasing map $f:[k]\to[n]$ such that $\orient(C_{f(i)},C_{f(j)},C_{f(l)})<0$ for
  \emph{all} $i<j<l$ in $[k]$.
\item For every finite set $T\subset\R^2$ let $\mathrm{xset}(T)=\{p_0:p\in T\}\subset\R$. From now on
  let $T$ also be $x$-injective. By Lemma~\ref{mm:lem:sort}(a) the set $\mathrm{xset}(T)$ has $\card T$
  elements, so there is a unique strictly increasing map $\mathrm{emb}(T):[\card T]\to\R$ with image
  $\mathrm{xset}(T)$ (Mathlib \lean{Finset.orderEmbOfFin}); $\mathrm{emb}(T)(i)$ is the $(i+1)$-st
  smallest abscissa of a point of $T$. For $i\in[\card T]$, $\mathrm{sortPt}(T)(i)$ is a point $p\in T$
  with $p_0=\mathrm{emb}(T)(i)$. Such a point exists by Lemma~\ref{mm:lem:sort}(b), and Lean selects one
  with the axiom of choice (\lean{Classical.choose}). By $x$-injectivity it is unique, but this is not
  used. (Parts (a) and (b) of Lemma~\ref{mm:lem:sort} concern only $\mathrm{xset}$ and $\mathrm{emb}$,
  so these forward references are not circular.)
\item If, in addition, $T$ is in general position, $\mathrm{cfg}(T)$ is the configuration of $\card T$ points
  $i\mapsto\mathrm{sortPt}(T)(i)$. Its two defining properties are Lemma~\ref{mm:lem:sort}(e) and~(h).
\end{enumerate}
\end{definition}

\begin{remark}
Caps in Definition~\ref{mm:def:config}(b) are \emph{all-triples} caps: every increasing triple of the
selected points is clockwise. The written sources \dirn{MIRRORETV/README.md} and
\dirn{MIRRORETV/PROOFS.md} define caps by consecutive triples; for points with increasing abscissae the
two notions agree (\dirn{MIRRORETV/PROOFS.md}, Lemma~2.1(b)), but only the all-triples notion occurs in
this section. For $k\le2$ the condition on triples is vacuous, so a configuration of $n$ points has a
$k$-cap if and only if $k\le n$.
\end{remark}

\begin{lemma}[sorting; \leanref{MirrorOdd.xset_card}{MIRRORETV/lean/MirrorMain.lean:1551}, \leanref{MirrorOdd.emb_mem}{MIRRORETV/lean/MirrorMain.lean:1555}, \leanref{MirrorOdd.sortPt_mem}{MIRRORETV/lean/MirrorMain.lean:1560}, \leanref{MirrorOdd.sortPt_x}{MIRRORETV/lean/MirrorMain.lean:1563}, \leanref{MirrorOdd.sortPt_strictMono}{MIRRORETV/lean/MirrorMain.lean:1566}, \leanref{MirrorOdd.sortPt_injective}{MIRRORETV/lean/MirrorMain.lean:1572}, \leanref{MirrorOdd.sortPt_image}{MIRRORETV/lean/MirrorMain.lean:1576}]\label{mm:lem:sort}
Let $T\subset\R^2$ be finite and $x$-injective, and write $s_i=\mathrm{sortPt}(T)(i)$ for
$i\in[\card T]$.
\begin{enumerate}[label=(\alph*)]
\item $\card{\mathrm{xset}(T)}=\card T$ (\lean{xset_card}).
\item For every $i\in[\card T]$ there is $p\in T$ with $p_0=\mathrm{emb}(T)(i)$ (\lean{emb_mem}).
\item $s_i\in T$ for every $i$ (\lean{sortPt_mem}).
\item $(s_i)_0=\mathrm{emb}(T)(i)$ for every $i$ (\lean{sortPt_x}).
\item $i\mapsto(s_i)_0$ is strictly increasing on $[\card T]$ (\lean{sortPt_strictMono}).
\item $i\mapsto s_i$ is injective (\lean{sortPt_injective}).
\item $\{s_i:i\in[\card T]\}=T$ (\lean{sortPt_image}).
\item (the field \lean{generalPosition} of \lean{cfg}, \dirn{MIRRORETV/lean/MirrorMain.lean:1587})
  If $T$ is in general position, then $\orient(s_i,s_j,s_k)\ne0$ for all $i<j<k$ in $[\card T]$.
\end{enumerate}
Consequently, if $T$ is in general position, $\mathrm{cfg}(T)$ is a configuration of $\card T$ points with
$\mathrm{pts}(\mathrm{cfg}(T))=T$.
\end{lemma}

\begin{proof}
(a) The map $p\mapsto p_0$ is injective on $T$, so its image $\mathrm{xset}(T)$ has $\card T$ elements.
(b) $\mathrm{emb}(T)(i)$ lies in the image $\mathrm{xset}(T)$ of $\mathrm{emb}(T)$, so it equals $p_0$ for
some $p\in T$.
(c), (d) These are the two properties of the point chosen in Definition~\ref{mm:def:config}(c).
(e) $\mathrm{emb}(T)$ is strictly increasing; use (d).
(f) If $s_i=s_j$, then $(s_i)_0=(s_j)_0$, so $i=j$ by (e).
(g) By (c) the set $I=\{s_i:i\in[\card T]\}$ is contained in $T$, and by (f) it has $\card T$ elements. A
subset of the finite set $T$ with $\card T$ elements is $T$ itself.
(h) By (f) the points $s_i,s_j,s_k$ are pairwise distinct, and by (c) they lie in $T$. Since $T$ is in
general position, $\orient(s_i,s_j,s_k)\ne0$ by Lemma~\ref{fd:lem:collinear}.
Finally, (e) and (h) are the two conditions of Definition~\ref{mm:def:config}(a) for $i\mapsto s_i$, and
(g) says that the point set is $T$.
\end{proof}

\begin{lemma}[two facts on convex polygons; \leanref{MirrorOdd.hasConvexNGon_of_le}{MIRRORETV/lean/MirrorMain.lean:2182}, \leanref{MirrorOdd.hasConvexNGon_affineMap}{MIRRORETV/lean/MirrorMain.lean:1859}]\label{mm:lem:ngon}
Let $S\subset\R^2$ be finite.
\begin{enumerate}[label=(\alph*)]
\item (\lean{hasConvexNGon_of_le}) If $S$ has a convex $m$-gon and $n\le m$, then $S$ has a convex
  $n$-gon.
\item (\lean{hasConvexNGon_affineMap}) If $S$ has a convex $n$-gon and $f:\R^2\to\R^2$ is an injective
  affine map, then $f(S)$ has a convex $n$-gon.
\end{enumerate}
\end{lemma}

\begin{proof}
(a) This is Lemma~\ref{fd:lem:mono}(d). Directly: take $T\subseteq S$ with $\card T=m$ in convex
position and choose $T'\subseteq T$ with $\card{T'}=n$. Then $T'\subseteq S$, and $T'$ is in convex
position by Lemma~\ref{fd:lem:mono}(a).

(b) Take $T\subseteq S$ with $\card T=n$ in convex position. Then $f(T)\subseteq f(S)$ and
$\card{f(T)}=n$ because $f$ is injective. An affine map preserves affine combinations:
$f(\sum_i\lambda_iy_i)=\sum_i\lambda_if(y_i)$ whenever $\sum_i\lambda_i=1$. Since the convex hull of a set
consists of the convex combinations of its finite subsets, $f(\conv Y)=\conv f(Y)$ for every
$Y\subseteq\R^2$. Let $q\in f(T)$, say $q=f(p)$ with $p\in T$. By injectivity,
$f(T)\setminus\{f(p)\}=f(T\setminus\{p\})$. If $f(p)\in\conv f(T\setminus\{p\})=f(\conv(T\setminus\{p\}))$,
then $f(p)=f(x)$ for some $x\in\conv(T\setminus\{p\})$, so $p=x\in\conv(T\setminus\{p\})$ by injectivity,
which contradicts the convex position of $T$ (Lemma~\ref{fd:lem:cp-hull}). Hence no point of $f(T)$ lies
in the convex hull of the other points of $f(T)$, and $f(T)$ is in convex position by
Lemma~\ref{fd:lem:cp-hull}. (This is also contained in Lemma~\ref{fd:lem:affine-cp}.)
\end{proof}

\subsection{The mirror lemma}

\begin{lemma}[cap facts; \leanref{MirrorOdd.cap_below_lines}{MIRRORETV/lean/MirrorMain.lean:2204}]\label{mm:lem:slope}
Both parts concern points that lie on or above the $x$-axis.
\begin{enumerate}[label=(\alph*)]
\item (\lean{cap_below_lines}) Let $m\in\N$ and let $c:[m]\to\R^2$ be such that $i\mapsto(c_i)_0$ is
  strictly increasing, $(c_i)_1\ge0$ for all $i\in[m]$, and $\orient(c_i,c_j,c_l)<0$ for all $i<j<l$ in
  $[m]$. Then for all $i,j,l\in[m]$ with $i<j$ and $(c_l)_1>0$,
  \[ \orient(c_i,c_j,\operatorname{mir}(c_l))<0. \]
\item (Written step. It is not a separate Lean declaration: the same argument is carried out inline in
  the proof of \lean{mirror_odd_general}, \dirn{MIRRORETV/lean/MirrorMain.lean:2542}--2557, and as the
  local claim \lean{interior} in \lean{mirror_five}, line~2060, and in \lean{mirror_seven_of},
  line~1716.) Let $p,a,q\in\R^2$ with $p_0<a_0<q_0$, $p_1\ge0$, $q_1\ge0$ and $a_1=0$. Then
  \[ \orient(p,a,q)=(a_0-p_0)\,q_1+p_1\,(q_0-a_0)\ \ge\ 0 . \]
  In particular $\orient(p,a,q)<0$ is impossible: a point on the axis cannot be the middle point of a
  clockwise triple with increasing abscissae whose other two points lie on or above the axis.
\end{enumerate}
\end{lemma}

In words, (a) says that the mirror image of a cap point lying strictly above the axis lies strictly to the
right of every directed chord $c_i\to c_j$ of the cap ($i<j$), that is, strictly below the line through
$c_i$ and $c_j$.

\begin{proof}
(a) Write $x_k=(c_k)_0$ and $y_k=(c_k)_1$. By Lemma~\ref{mm:lem:orient-mir}(e),
\[ \orient(c_i,c_j,\operatorname{mir}c_l)=\orient(c_i,c_j,c_l)-2y_l(x_j-x_i), \tag{$*$} \]
and $2y_l(x_j-x_i)>0$ because $y_l>0$ and $x_j>x_i$. We distinguish the position of $l$.
\begin{itemize}
\item $l<i$: $\orient(c_i,c_j,c_l)=\orient(c_l,c_i,c_j)$ by Lemma~\ref{mm:lem:orient-mir}(d), and this is
  negative by the cap hypothesis for $l<i<j$. So $(*)$ is negative.
\item $l=i$ or $l=j$: $\orient(c_i,c_j,c_l)=0$ because two arguments coincide, so $(*)$ equals
  $-2y_l(x_j-x_i)<0$.
\item $i<l<j$: expanding directly,
  \[ \orient(c_i,c_j,\operatorname{mir}c_l)=-(x_j-x_i)\,y_l-y_i\,(x_j-x_l)-y_j\,(x_l-x_i). \]
  The first term is negative, and the other two are $\le0$ because $y_i,y_j\ge0$ and $x_i<x_l<x_j$.
  (The cap hypothesis is not needed in this case.)
\item $l>j$: $\orient(c_i,c_j,c_l)<0$ by the cap hypothesis for $i<j<l$, so $(*)$ is negative.
\end{itemize}
(b) Expanding with $a_1=0$:
$\orient(p,a,q)=(a_0-p_0)(q_1-p_1)+p_1(q_0-p_0)=(a_0-p_0)q_1+p_1(q_0-a_0)$, and both products are $\ge0$.
\end{proof}

\begin{theorem}[mirror lemma, no point on the axis; \leanref{MirrorOdd.cap_mirror_polygon}{MIRRORETV/lean/MirrorMain.lean:2233}]\label{mm:thm:mirror-polygon}
Let $m\in\N$ with $2\le m$, and let $c:[m]\to\R^2$ be such that $i\mapsto(c_i)_0$ is strictly increasing,
$(c_i)_1>0$ for all $i\in[m]$, and $\orient(c_i,c_j,c_l)<0$ for all $i<j<l$ in $[m]$. Let $P\subset\R^2$
be a finite set with $c_i\in P$ and $\operatorname{mir}(c_i)\in P$ for all $i\in[m]$. Then $P$ has a
convex $2m$-gon.
\end{theorem}

\begin{proof}
We follow the Lean proof. Write $x_i=(c_i)_0$ and $y_i=(c_i)_1>0$. Define $q:[2m]\to\R^2$ by
\[ q_r=\operatorname{mir}(c_r)\quad(0\le r\le m-1),\qquad q_r=c_{2m-1-r}\quad(m\le r\le 2m-1). \]
Thus $q$ lists $\operatorname{mir}c_0,\dots,\operatorname{mir}c_{m-1}$ (the reflected cap, from left to
right) followed by $c_{m-1},\dots,c_0$ (the cap, from right to left). For $m\le r\le 2m-1$ the index
$2m-1-r$ lies in $[m]$, and $r\mapsto2m-1-r$ is strictly decreasing. Since all $y_i>0$, the hypotheses
of Lemma~\ref{mm:lem:slope}(a) hold for $c$.

\emph{Claim~1: $\orient(q_r,q_s,q_t)>0$ for all $r<s<t$ in $[2m]$.} We distinguish how many of $r,s,t$
are smaller than $m$.
\begin{itemize}
\item Case~1, $t<m$. By Lemma~\ref{mm:lem:orient-mir}(b),
  $\orient(\operatorname{mir}c_r,\operatorname{mir}c_s,\operatorname{mir}c_t)=-\orient(c_r,c_s,c_t)$, which
  is positive by the cap hypothesis for $r<s<t$.
\item Case~2, $s<m\le t$. Put $u=2m-1-t\in[m]$, so $q_t=c_u$. By Lemma~\ref{mm:lem:orient-mir}(c),
  $\orient(\operatorname{mir}c_r,\operatorname{mir}c_s,c_u)=-\orient(c_r,c_s,\operatorname{mir}c_u)$, and
  $\orient(c_r,c_s,\operatorname{mir}c_u)<0$ by Lemma~\ref{mm:lem:slope}(a) with $(i,j,l)=(r,s,u)$, since
  $r<s$ and $y_u>0$.
\item Case~3, $r<m\le s$. Put $u'=2m-1-s$ and $u=2m-1-t$. Then $u<u'$ (because $s<t$), $q_s=c_{u'}$ and
  $q_t=c_u$. By the first two identities of Lemma~\ref{mm:lem:orient-mir}(d),
  \[ \orient(\operatorname{mir}c_r,c_{u'},c_u)=\orient(c_{u'},c_u,\operatorname{mir}c_r)
     =-\orient(c_u,c_{u'},\operatorname{mir}c_r), \]
  and $\orient(c_u,c_{u'},\operatorname{mir}c_r)<0$ by Lemma~\ref{mm:lem:slope}(a) with $(i,j,l)=(u,u',r)$,
  since $u<u'$ and $y_r>0$.
\item Case~4, $m\le r$. Put $\alpha=2m-1-r$, $\beta=2m-1-s$ and $\delta=2m-1-t$; then
  $\delta<\beta<\alpha$, $q_r=c_\alpha$, $q_s=c_\beta$ and $q_t=c_\delta$. By the third identity of
  Lemma~\ref{mm:lem:orient-mir}(d), $\orient(c_\alpha,c_\beta,c_\delta)=-\orient(c_\delta,c_\beta,c_\alpha)$,
  which is positive by the cap hypothesis for $\delta<\beta<\alpha$.
\end{itemize}

\emph{Claim~2: $q$ is injective.} Let $r\ne s$ in $[2m]$. If $r,s<m$ and $q_r=q_s$, then comparing
abscissae gives $x_r=x_s$ (Lemma~\ref{mm:lem:orient-mir}(a)), so $r=s$ because $i\mapsto x_i$ is strictly
increasing; a contradiction. If $r<m\le s$, then $(q_r)_1=-y_r<0<y_{2m-1-s}=(q_s)_1$, so $q_r\ne q_s$; the
case $s<m\le r$ is symmetric. If $r,s\ge m$ and $q_r=q_s$, then $x_{2m-1-r}=x_{2m-1-s}$, so
$2m-1-r=2m-1-s$ and $r=s$; a contradiction.

\emph{Conclusion.} Every $q_r$ lies in $P$, being some $\operatorname{mir}(c_i)$ or some $c_i$. Apply
Lemma~\ref{fd:lem:pos-cycle} with $N=k=2m$ (note $2m\ge4$ because $m\ge2$), the injective map $p=q$
(Claim~2), $e$ the identity of $[2m]$ and the finite set $P$; its positivity hypothesis is Claim~1. Hence
$P$ has a convex $2m$-gon.
\end{proof}

\begin{remark}
(1) Claim~1 checks the sign of \emph{every} triple $q_r,q_s,q_t$ with $r<s<t$, that is, of every triple
of the list taken in list order, which is what Lemma~\ref{fd:lem:pos-cycle} requires. It would not suffice to check that the closed polygon
$q_0q_1\cdots q_{2m-1}$ turns the same way at every vertex: such a polygon can wind around twice, and then
a vertex can lie inside the convex hull of the others (\dirn{MIRRORETV/PROOFS.md}, remark after Lemma~3.1,
gives a five-point example).
(2) \dirn{MIRRORETV/PROOFS.md}, Lemma~3.1, gives a written argument by a different method (a region
argument) for a version with caps $C$ defined by consecutive triples: if the interior points of $C$ have
positive ordinates and its endpoints have nonnegative ordinates (positive if $C$ has only one point), then
$C\cup\operatorname{mir}(C)$ is in convex position. That argument is not used here; the Lean statements
are Theorems~\ref{mm:thm:mirror-polygon} and~\ref{mm:thm:axis-polygon}, and these are what we prove.
\end{remark}

\begin{theorem}[mirror lemma, one endpoint on the axis; \leanref{MirrorOdd.axis_left_polygon}{MIRRORETV/lean/MirrorMain.lean:2309}, \leanref{MirrorOdd.axis_right_polygon}{MIRRORETV/lean/MirrorMain.lean:2403}]\label{mm:thm:axis-polygon}
Let $m\in\N$ with $3\le m$, and let $c:[m]\to\R^2$ be such that $i\mapsto(c_i)_0$ is strictly increasing
and $\orient(c_i,c_j,c_l)<0$ for all $i<j<l$ in $[m]$. Let $P\subset\R^2$ be a finite set with $c_i\in P$
and $\operatorname{mir}(c_i)\in P$ for all $i\in[m]$.
\begin{enumerate}[label=(\alph*)]
\item (\lean{axis_left_polygon}) If $(c_0)_1=0$ and $(c_i)_1>0$ for every $i\in[m]$ with $0<i$, then $P$
  has a convex $(2m-1)$-gon.
\item (\lean{axis_right_polygon}) If $(c_{m-1})_1=0$ and $(c_i)_1>0$ for every $i\in[m]$ with $i<m-1$,
  then $P$ has a convex $(2m-1)$-gon.
\end{enumerate}
\end{theorem}

\begin{proof}
We follow the Lean proofs. Write $x_i=(c_i)_0$ and $y_i=(c_i)_1$. In both parts $y_i\ge0$ for all $i$,
so Lemma~\ref{mm:lem:slope}(a) applies to $c$; and $2m-1\ge5$.

(a) Here $\operatorname{mir}(c_0)=c_0$ because $(c_0)_1=0$. Define $q:[2m-1]\to\R^2$ by
\[ q_r=\operatorname{mir}(c_r)\quad(0\le r\le m-1),\qquad q_r=c_{2m-1-r}\quad(m\le r\le 2m-2). \]
So $q$ lists $\operatorname{mir}c_0=c_0,\operatorname{mir}c_1,\dots,\operatorname{mir}c_{m-1}$ and then
$c_{m-1},\dots,c_1$. For $m\le r\le 2m-2$ the index $2m-1-r$ lies in $\{1,\dots,m-1\}$, hence
$y_{2m-1-r}>0$.

\emph{Claim~1: $\orient(q_r,q_s,q_t)>0$ for all $r<s<t$ in $[2m-1]$.}
\begin{itemize}
\item Case~1 ($t<m$) and Case~4 ($m\le r$) are exactly as in the proof of
  Theorem~\ref{mm:thm:mirror-polygon}; they use only the cap hypothesis (in Case~4 the indices are
  $2m-1-r>2m-1-s>2m-1-t$).
\item Case~2, $s<m\le t$. With $u=2m-1-t\in\{1,\dots,m-1\}$ we have $q_t=c_u$ and $y_u>0$, and Case~2 of
  the proof of Theorem~\ref{mm:thm:mirror-polygon} applies verbatim.
\item Case~3, $r<m\le s$. Put $u'=2m-1-s$ and $u=2m-1-t$, so $1\le u<u'\le m-1$, $q_s=c_{u'}$, $q_t=c_u$.
  \begin{itemize}
  \item Case~3a, $r=0$. Then $q_r=\operatorname{mir}c_0=c_0$, and by Lemma~\ref{fd:lem:orient}(a)
    (exchange of the last two arguments) $\orient(c_0,c_{u'},c_u)=-\orient(c_0,c_u,c_{u'})$, which is
    positive by the cap hypothesis for $0<u<u'$.
  \item Case~3b, $r\ge1$. Then $y_r>0$, and Case~3 of the proof of Theorem~\ref{mm:thm:mirror-polygon}
    applies verbatim:
    $\orient(\operatorname{mir}c_r,c_{u'},c_u)=-\orient(c_u,c_{u'},\operatorname{mir}c_r)>0$.
  \end{itemize}
\end{itemize}

\emph{Claim~2: $q$ is injective.} Two indices $r\ne s$ that are both $<m$, or both $\ge m$, are separated by
the abscissae as in Theorem~\ref{mm:thm:mirror-polygon}. If $r<m\le s$, then $(q_s)_1=y_{2m-1-s}>0$
because $2m-1-s\ge1$, while $(q_r)_1=-y_r\le0$. The case $s<m\le r$ is symmetric.

\emph{Conclusion.} Every $q_r$ lies in $P$, and Lemma~\ref{fd:lem:pos-cycle} with $N=k=2m-1\ge4$, $p=q$
and $e$ the identity gives a convex $(2m-1)$-gon in $P$.

(b) Here $\operatorname{mir}(c_{m-1})=c_{m-1}$. Define $q:[2m-1]\to\R^2$ by
\[ q_r=\operatorname{mir}(c_r)\quad(0\le r\le m-1),\qquad q_r=c_{2m-2-r}\quad(m\le r\le 2m-2). \]
So $q$ lists $\operatorname{mir}c_0,\dots,\operatorname{mir}c_{m-2},\operatorname{mir}c_{m-1}=c_{m-1}$ and
then $c_{m-2},\dots,c_0$. For $m\le r\le 2m-2$ the index $2m-2-r$ lies in $\{0,\dots,m-2\}$, hence
$y_{2m-2-r}>0$.

\emph{Claim~1.}
\begin{itemize}
\item Cases~1 and~4 are as in Theorem~\ref{mm:thm:mirror-polygon}, with $2m-2-\cdot$ in place of
  $2m-1-\cdot$.
\item Case~2, $s<m\le t$. With $u=2m-2-t\le m-2$ we have $q_t=c_u$ and $y_u>0$, and Case~2 of
  Theorem~\ref{mm:thm:mirror-polygon} applies.
\item Case~3, $r<m\le s$. Put $u'=2m-2-s$ and $u=2m-2-t$, so $0\le u<u'\le m-2$.
  \begin{itemize}
  \item Case~3a, $r=m-1$. Then $q_r=c_{m-1}$, and by the third identity of
    Lemma~\ref{mm:lem:orient-mir}(d), $\orient(c_{m-1},c_{u'},c_u)=-\orient(c_u,c_{u'},c_{m-1})$, which is
    positive by the cap hypothesis for $u<u'<m-1$.
  \item Case~3b, $r<m-1$. Then $y_r>0$, and Case~3 of Theorem~\ref{mm:thm:mirror-polygon} applies.
  \end{itemize}
\end{itemize}

\emph{Claim~2.} Two indices $r\ne s$ that are both $<m$, or both $\ge m$, are separated by the abscissae
as in (a). If $r<m\le s$, then $(q_s)_1=y_{2m-2-s}>0$ because $2m-2-s\le m-2$, while
$(q_r)_1=-y_r\le0$. The case $s<m\le r$ is symmetric.

\emph{Conclusion} as in (a).
\end{proof}

\begin{remark}
The hypothesis $m\ge3$ makes the list long enough for Lemma~\ref{fd:lem:pos-cycle}, which needs at least
four points. In (a) the point $c_0$ occurs in the list only once, as $\operatorname{mir}c_0=c_0$; in (b)
the same holds for $c_{m-1}$. This is why one point is lost compared with Theorem~\ref{mm:thm:mirror-polygon}.
\end{remark}

\subsection{The generated small polygons}

Lines 22--1506 of \dirn{MIRRORETV/lean/MirrorMain.lean} contain five helper lemmas $\mathrm{cyc}_k$
($k=4,\dots,8$) and nine lemmas on small mirror polygons. In Lean each polygon lemma is proved directly:
the pairwise distinctness of its $k$ vertices by comparing one coordinate, and each of the $\binom k3$
orientation inequalities, after expanding $\orient$ and $\operatorname{mir}$, by linear arithmetic
(\lean{linarith}) from the hypotheses (with the cap orientations expanded) and auxiliary products
$(c_i)_1\,(u_0-v_0)>0$, where $u$ and $v$ are two of the lemma's points with $v_0<u_0$ (in the axis
lemmas one of them can be the axis point $a$), with certificates found offline; then $\mathrm{cyc}_k$ is
applied. These Lean proofs do not use the
general lemmas, which come later in the file. Our written proofs are different: we derive the nine lemmas
from Theorems~\ref{mm:thm:mirror-polygon} and~\ref{mm:thm:axis-polygon}.

\begin{lemma}[positive cycles of length $4$ to $8$; \leanref{MirrorOdd.cyc4}{MIRRORETV/lean/MirrorMain.lean:995}, \leanref{MirrorOdd.cyc5}{MIRRORETV/lean/MirrorMain.lean:1026}, \leanref{MirrorOdd.cyc6}{MIRRORETV/lean/MirrorMain.lean:1068}, \leanref{MirrorOdd.cyc7}{MIRRORETV/lean/MirrorMain.lean:22}, \leanref{MirrorOdd.cyc8}{MIRRORETV/lean/MirrorMain.lean:103}]\label{mm:lem:cyc}
Let $k\in\{4,5,6,7,8\}$, let $q_0,\dots,q_{k-1}\in\R^2$, and let $P\subset\R^2$ be finite. Suppose that
$q_r\in P$ for all $r\in[k]$, that $q_r\ne q_s$ for all $r<s$ in $[k]$, and that
$\orient(q_r,q_s,q_t)>0$ for all $r<s<t$ in $[k]$. Then $P$ has a convex $k$-gon. (The Lean lemma
$\mathrm{cyc}_k$ lists these $k$ memberships, $\binom k2$ inequalities $q_r\ne q_s$ and $\binom k3$
orientation inequalities as separate hypotheses.)
\end{lemma}

\begin{proof}
The map $r\mapsto q_r$ on $[k]$ is injective by the distinctness hypotheses. Lemma~\ref{fd:lem:pos-cycle}
with $N=k\ge4$, $p=q$, $e$ the identity of $[k]$ and the set $P$ gives the claim; its remaining hypotheses
are exactly the memberships and the orientation inequalities. This is also the Lean proof.
\end{proof}

\begin{proposition}[generated small mirror polygons; \leanref{MirrorOdd.cap_mirror_quadrilateral}{MIRRORETV/lean/MirrorMain.lean:1126}, \leanref{MirrorOdd.cap_mirror_pentagon}{MIRRORETV/lean/MirrorMain.lean:1282}, \leanref{MirrorOdd.cap_mirror_hexagon}{MIRRORETV/lean/MirrorMain.lean:1168}, \leanref{MirrorOdd.cap_mirror_heptagon}{MIRRORETV/lean/MirrorMain.lean:212}, \leanref{MirrorOdd.cap_mirror_octagon}{MIRRORETV/lean/MirrorMain.lean:392}, \leanref{MirrorOdd.axis_left_pentagon}{MIRRORETV/lean/MirrorMain.lean:1358}, \leanref{MirrorOdd.axis_right_pentagon}{MIRRORETV/lean/MirrorMain.lean:1433}, \leanref{MirrorOdd.axis_left_heptagon}{MIRRORETV/lean/MirrorMain.lean:639}, \leanref{MirrorOdd.axis_right_heptagon}{MIRRORETV/lean/MirrorMain.lean:816}]\label{mm:prop:small-polygons}
Let $a,c_0,c_1,c_2,c_3\in\R^2$ and let $P\subset\R^2$ be finite. Write $x_i=(c_i)_0$ and $y_i=(c_i)_1$.
\begin{enumerate}[label=(\alph*)]
\item (\lean{cap_mirror_quadrilateral}) If $x_0<x_1$, $y_0>0$, $y_1>0$, and
  $\operatorname{mir}c_0,\operatorname{mir}c_1,c_1,c_0\in P$, then $P$ has a convex $4$-gon.
\item (\lean{cap_mirror_pentagon}) If $x_0<x_1<x_2$, $y_0,y_1,y_2>0$, $\orient(c_0,c_1,c_2)<0$, and
  $\operatorname{mir}c_0,\operatorname{mir}c_1,\operatorname{mir}c_2,c_2,c_1\in P$, then $P$ has a convex
  $5$-gon.
\item (\lean{cap_mirror_hexagon}) If $x_0<x_1<x_2$, $y_0,y_1,y_2>0$, $\orient(c_0,c_1,c_2)<0$, and
  $\operatorname{mir}c_0,\operatorname{mir}c_1,\operatorname{mir}c_2,c_2,c_1,c_0\in P$, then $P$ has a
  convex $6$-gon.
\item (\lean{cap_mirror_heptagon}) If $x_0<x_1<x_2<x_3$, $y_0,y_1,y_2,y_3>0$, the four orientations
  $\orient(c_0,c_1,c_2)$, $\orient(c_0,c_1,c_3)$, $\orient(c_0,c_2,c_3)$, $\orient(c_1,c_2,c_3)$ are $<0$,
  and $c_0,\operatorname{mir}c_1,\operatorname{mir}c_2,\operatorname{mir}c_3,c_3,c_2,c_1\in P$, then $P$
  has a convex $7$-gon.
\item (\lean{cap_mirror_octagon}) Under the hypotheses on $c_0,\dots,c_3$ of (d), if
  $\operatorname{mir}c_0,\operatorname{mir}c_1,\operatorname{mir}c_2,\operatorname{mir}c_3,c_3,c_2,c_1,c_0\in P$,
  then $P$ has a convex $8$-gon.
\item (\lean{axis_left_pentagon}) If $a_1=0$, $a_0<x_0<x_1$, $y_0,y_1>0$, $\orient(a,c_0,c_1)<0$, and
  $a,\operatorname{mir}c_0,\operatorname{mir}c_1,c_1,c_0\in P$, then $P$ has a convex $5$-gon.
\item (\lean{axis_right_pentagon}) If $a_1=0$, $x_0<x_1<a_0$, $y_0,y_1>0$, $\orient(c_0,c_1,a)<0$, and
  $\operatorname{mir}c_0,\operatorname{mir}c_1,a,c_1,c_0\in P$, then $P$ has a convex $5$-gon.
\item (\lean{axis_left_heptagon}) If $a_1=0$, $a_0<x_0<x_1<x_2$, $y_0,y_1,y_2>0$, the four orientations
  $\orient(a,c_0,c_1)$, $\orient(a,c_0,c_2)$, $\orient(a,c_1,c_2)$, $\orient(c_0,c_1,c_2)$ are $<0$, and
  $a,\operatorname{mir}c_0,\operatorname{mir}c_1,\operatorname{mir}c_2,c_2,c_1,c_0\in P$, then $P$ has a
  convex $7$-gon.
\item (\lean{axis_right_heptagon}) If $a_1=0$, $x_0<x_1<x_2<a_0$, $y_0,y_1,y_2>0$, the four orientations
  $\orient(c_0,c_1,c_2)$, $\orient(c_0,c_1,a)$, $\orient(c_0,c_2,a)$, $\orient(c_1,c_2,a)$ are $<0$, and
  $\operatorname{mir}c_0,\operatorname{mir}c_1,\operatorname{mir}c_2,a,c_2,c_1,c_0\in P$, then $P$ has a
  convex $7$-gon.
\end{enumerate}
In each item the memberships are listed in the vertex order that the Lean proof passes to
$\mathrm{cyc}_k$ (Lemma~\ref{mm:lem:cyc}).
\end{proposition}

\begin{proof}
We use the general theorems; this differs from the generated Lean proofs described above.

(a), (c), (e): apply Theorem~\ref{mm:thm:mirror-polygon} with $m=2,3,4$ to $c=(c_0,\dots,c_{m-1})$ and to
$P$. The abscissae increase strictly because consecutive ones do, and all $y_i>0$. The cap hypothesis
concerns the triples $i<j<l$ in $[m]$: there are none for $m=2$, only $(0,1,2)$ for $m=3$, and
$(0,1,2)$, $(0,1,3)$, $(0,2,3)$, $(1,2,3)$ for $m=4$; these are exactly the listed orientation hypotheses.
$P$ contains all $c_i$ and $\operatorname{mir}c_i$. The conclusion is a convex $2m$-gon, that is, a
$4$-, $6$- or $8$-gon.

(f), (h): note that $\operatorname{mir}(a)=(a_0,-0)=a$. Apply Theorem~\ref{mm:thm:axis-polygon}(a) with
$m=3$ to $(a,c_0,c_1)$, respectively with $m=4$ to $(a,c_0,c_1,c_2)$. The abscissae increase strictly, the
first point $a$ has ordinate $0$ and the others have positive ordinates. For $m=3$ the only triple is
$(a,c_0,c_1)$; for $m=4$ the triples $(0,1,2)$, $(0,1,3)$, $(0,2,3)$, $(1,2,3)$ of $(a,c_0,c_1,c_2)$ are
$(a,c_0,c_1)$, $(a,c_0,c_2)$, $(a,c_1,c_2)$, $(c_0,c_1,c_2)$. These are the listed orientation hypotheses.
$P$ contains the listed points and their mirror images (recall $\operatorname{mir}a=a$). The conclusion is
a convex $(2m-1)$-gon, that is, a $5$- or $7$-gon.

(g), (i): the same with Theorem~\ref{mm:thm:axis-polygon}(b), applied to $(c_0,c_1,a)$, respectively
$(c_0,c_1,c_2,a)$, whose last point $a$ has ordinate $0$. For $m=4$ the four triples are
$(c_0,c_1,c_2)$, $(c_0,c_1,a)$, $(c_0,c_2,a)$, $(c_1,c_2,a)$.

(b): let $Q=\{\operatorname{mir}c_0,\operatorname{mir}c_1,\operatorname{mir}c_2,c_2,c_1,c_0\}$, a set with
at most $6$ elements. By (c) applied to the finite set $Q$ in place of $P$, there is $T\subseteq Q$ with
$\card T=6$ in convex position. As $\card Q\le6$, we get $T=Q$: the six listed points are pairwise distinct
and $Q$ is in convex position. Then $Q\setminus\{c_0\}=\{\operatorname{mir}c_0,\operatorname{mir}c_1,
\operatorname{mir}c_2,c_2,c_1\}$ has $5$ elements, is in convex position (Lemma~\ref{fd:lem:mono}(a)), and
is contained in $P$ by hypothesis. So $P$ has a convex $5$-gon.

(d): in the same way, (e) applied to the set $Q$ of its eight points shows that these eight points are
pairwise distinct and in convex position. The set $Q\setminus\{\operatorname{mir}c_0\}=\{c_0,
\operatorname{mir}c_1,\operatorname{mir}c_2,\operatorname{mir}c_3,c_3,c_2,c_1\}$ has $7$ elements, is in
convex position and is contained in $P$.
\end{proof}

In Lean, items (a), (b), (c), (f), (g) are used in the proofs of Theorem~\ref{mm:thm:small}, and items
(d), (e), (h), (i) in those of the bounds for $k=7,8$ (Theorems~\ref{mb:thm:seven}
and~\ref{mb:thm:eight}).

\subsection{Counting and distinct abscissae}

\begin{lemma}[counting; \leanref{MirrorOdd.card_split}{MIRRORETV/lean/MirrorMain.lean:1601}, \leanref{MirrorOdd.card_A_le}{MIRRORETV/lean/MirrorMain.lean:1636}]\label{mm:lem:count}
Let $S\subset\R^2$ be finite.
\begin{enumerate}[label=(\alph*)]
\item (\lean{card_split}) If $\operatorname{mir}(p)\in S$ for every $p\in S$, then
  $\card S=2\card{R(S)}+\card{A(S)}$. (General position is not assumed.)
\item (\lean{card_A_le}) If $S$ is in general position, then $\card{A(S)}\le2$.
\end{enumerate}
\end{lemma}

\begin{proof}
(a) Let $N=\{p\in S:p_1<0\}$. By trichotomy of $p_1$, $S$ is the disjoint union of $R(S)$, $A(S)$ and $N$.
We show $\operatorname{mir}(R(S))=N$. If $q\in R(S)$, then $\operatorname{mir}q\in S$ by hypothesis and
$(\operatorname{mir}q)_1=-q_1<0$, so $\operatorname{mir}q\in N$. Conversely, if $p\in N$, then
$\operatorname{mir}p\in S$ and $(\operatorname{mir}p)_1=-p_1>0$, so $\operatorname{mir}p\in R(S)$, and
$p=\operatorname{mir}(\operatorname{mir}p)$ by Lemma~\ref{mm:lem:orient-mir}(a). Since $\operatorname{mir}$
is injective (it is its own inverse), $\card N=\card{R(S)}$, and therefore
$\card S=\card{R(S)}+\card{A(S)}+\card N=2\card{R(S)}+\card{A(S)}$.

(b) Suppose that $A(S)$ contains three pairwise distinct points $a,b,d$. Since $a_1=b_1=d_1=0$,
$\orient(a,b,d)=(b_0-a_0)(d_1-a_1)-(b_1-a_1)(d_0-a_0)=0$, which contradicts general position
(Lemma~\ref{fd:lem:collinear}). Hence $\card{A(S)}\le2$.
\end{proof}

\begin{lemma}[distinct abscissae; \leanref{MirrorOdd.injOn_x}{MIRRORETV/lean/MirrorMain.lean:1647}]\label{mm:lem:injx}
Let $S\subset\R^2$ be finite and in general position, with $\operatorname{mir}(p)\in S$ for every $p\in S$.
Let $T\subseteq S$ be a set such that $p_1\ge0$ for every $p\in T$, and such that $p=q$ whenever $p,q\in T$
and $p_1=q_1=0$ (at most one point of $T$ lies on the axis). Then $T$ is $x$-injective.
\end{lemma}

\begin{proof}
First, if $p,q\in T$, $p_1>0$ and $p_0=q_0$, then $p=q$. Indeed, suppose $p\ne q$. The points $p$, $q$
and $\operatorname{mir}p$ lie in $S$ (as $T\subseteq S$ and $S$ is $\operatorname{mir}$-symmetric) and are
pairwise distinct: $p\ne\operatorname{mir}p$ because $p_1>0>-p_1$, and $q\ne\operatorname{mir}p$ because
$q_1\ge0>-p_1$. But
$\orient(p,q,\operatorname{mir}p)=(q_0-p_0)(-p_1-p_1)-(q_1-p_1)(p_0-p_0)=0$ because $q_0=p_0$, which
contradicts general position (Lemma~\ref{fd:lem:collinear}).

Now let $p,q\in T$ with $p_0=q_0$. If $p_1>0$, then $p=q$ by the first step. Otherwise $p_1=0$. If
$q_1>0$, the first step with the roles of $p$ and $q$ exchanged gives $q=p$. Otherwise $q_1=0$ as well,
and $p=q$ by the hypothesis on axis points.
\end{proof}

\begin{lemma}[caps in the closed upper half; written step]\label{mm:lem:upper-caps}
Let $S\subset\R^2$ be finite and in general position, with $\operatorname{mir}(p)\in S$ for every $p\in S$.
Let $T\subseteq S$ satisfy $p_1\ge0$ for every $p\in T$, and $p=q$ whenever $p,q\in T$ and $p_1=q_1=0$.
Then $T$ is $x$-injective and in general position, so $C=\mathrm{cfg}(T)$ is a configuration of
$\card T$ points with $\mathrm{pts}(C)=T$. Let $m\ge2$ and suppose that $C$ has an $m$-cap.
\begin{enumerate}[label=(\alph*)]
\item If $p_1>0$ for every $p\in T$, then $S$ has a convex $2m$-gon.
\item If $m\ge3$, then $S$ has a convex $(2m-1)$-gon.
\end{enumerate}
(This lemma isolates the cap argument that the Lean proofs of Theorems~\ref{mm:thm:odd}
and~\ref{mm:thm:even} and of parts (b) and~(c) of Theorem~\ref{mm:thm:small} carry out inside their local
claims \lean{keyT} (odd case and $k=5$) and \lean{hR} (even case and $k=6$).)
\end{lemma}

\begin{proof}
$T$ is $x$-injective by Lemma~\ref{mm:lem:injx} and in general position by Lemma~\ref{fd:lem:mono}(b); the
statements about $C$ are Lemma~\ref{mm:lem:sort}. Let $f:[m]\to[\card T]$ be strictly increasing with
$\orient(C_{f(i)},C_{f(i')},C_{f(i'')})<0$ for all $i<i'<i''$ in $[m]$, and put $c_i=C_{f(i)}$ for
$i\in[m]$. Then $i\mapsto(c_i)_0$ is strictly increasing (a composition of strictly increasing maps);
$\orient(c_i,c_{i'},c_{i''})<0$ for all $i<i'<i''$; $c_i\in T\subseteq S$, hence
$\operatorname{mir}(c_i)\in S$; and $(c_i)_1\ge0$.

(a) Here $(c_i)_1>0$ for all $i$, because $c_i\in T$. Theorem~\ref{mm:thm:mirror-polygon} with $P=S$ gives
a convex $2m$-gon in $S$.

(b) If $(c_i)_1>0$ for all $i$, Theorem~\ref{mm:thm:mirror-polygon} gives a convex $2m$-gon in $S$, hence a
convex $(2m-1)$-gon by Lemma~\ref{mm:lem:ngon}(a). Otherwise there is $i_0$ with $(c_{i_0})_1\le0$, hence
$(c_{i_0})_1=0$. For $i\ne i_0$ we have $(c_i)_1>0$: otherwise $(c_i)_1=0$, the hypothesis on axis points
gives $c_i=c_{i_0}$, so $(c_i)_0=(c_{i_0})_0$ and $i=i_0$. There are three cases.
\begin{itemize}
\item $i_0=0$: Theorem~\ref{mm:thm:axis-polygon}(a) with $P=S$ gives a convex $(2m-1)$-gon in $S$.
\item $i_0=m-1$: Theorem~\ref{mm:thm:axis-polygon}(b) with $P=S$ gives a convex $(2m-1)$-gon in $S$.
\item $0<i_0<m-1$: put $p=c_{i_0-1}$, $a=c_{i_0}$ and $q=c_{i_0+1}$. Then $p_0<a_0<q_0$, $p_1,q_1\ge0$ and
  $a_1=0$, so $\orient(p,a,q)\ge0$ by Lemma~\ref{mm:lem:slope}(b); but $\orient(p,a,q)<0$ by the cap
  property for $i_0-1<i_0<i_0+1$. So this case does not occur. \qedhere
\end{itemize}
\end{proof}

\subsection{The mirror reduction}

\begin{theorem}[Theorem A, odd case; \leanref{MirrorOdd.mirror_odd_general}{MIRRORETV/lean/MirrorMain.lean:2498}]\label{mm:thm:odd}
Let $j,M\in\N$ with $2\le j$, and assume the following parametric hypothesis on $M$, which is part of the
statement (Remark~\ref{mm:rem:hm}):
\begin{itemize}
\item[(hM)] for every $n\in\N$ and every configuration $C$ of $n$ points: if $C$ has no $(j+1)$-cap and
  $\mathrm{pts}(C)$ has no convex $(2j+1)$-gon, then $n\le M$.
\end{itemize}
Let $S\subset\R^2$ be finite and in general position, with $\operatorname{mir}(p)\in S$ for every $p\in S$,
and suppose that $S$ has no convex $(2j+1)$-gon. Then $\card S\le2M$.
\end{theorem}

\begin{proof}
We follow the Lean proof.

\emph{Step~1.} Let $T\subseteq S$ be such that $p_1\ge0$ for all $p\in T$, and $p=q$ whenever $p,q\in T$
and $p_1=q_1=0$. Then $\card T\le M$. Indeed, let $C=\mathrm{cfg}(T)$, a configuration of $\card T$ points
with $\mathrm{pts}(C)=T$ (Lemma~\ref{mm:lem:upper-caps}). If $C$ had a $(j+1)$-cap,
Lemma~\ref{mm:lem:upper-caps}(b) with $m=j+1\ge3$ would give a convex $(2(j+1)-1)$-gon, that is, a convex
$(2j+1)$-gon, in $S$; so $C$ has no $(j+1)$-cap. If $\mathrm{pts}(C)=T$ had a convex $(2j+1)$-gon, so would
$S\supseteq T$ (Lemma~\ref{fd:lem:mono}(c)). Hypothesis (hM) with $n=\card T$ gives $\card T\le M$.

\emph{Step~2.} The set $T=R(S)$ satisfies the assumptions of Step~1: all its points have positive
ordinate, so none of them lies on the axis. Hence $\card{R(S)}\le M$.

\emph{Step~3.} If $A(S)=\emptyset$, then $\card S=2\card{R(S)}\le2M$ by Lemma~\ref{mm:lem:count}(a).
Otherwise pick $a\in A(S)$ and let $T=R(S)\cup\{a\}$. Then $a\notin R(S)$ (as $a_1=0$), $T\subseteq S$,
$p_1\ge0$ on $T$, and $a$ is the only point of $T$ on the axis. Step~1 gives
$\card{R(S)}+1=\card T\le M$. By Lemma~\ref{mm:lem:count}(a) and~(b),
$\card S=2\card{R(S)}+\card{A(S)}\le2\card{R(S)}+2\le2M$.
\end{proof}

\begin{theorem}[Theorem A, even case; \leanref{MirrorOdd.mirror_even_general}{MIRRORETV/lean/MirrorMain.lean:2599}]\label{mm:thm:even}
Let $j,M\in\N$ with $2\le j$, and assume the following parametric hypothesis on $M$, which is part of the
statement (Remark~\ref{mm:rem:hm}):
\begin{itemize}
\item[(hM)] for every $n\in\N$ and every configuration $C$ of $n$ points: if $C$ has no $j$-cap and
  $\mathrm{pts}(C)$ has no convex $2j$-gon, then $n\le M$.
\end{itemize}
Let $S\subset\R^2$ be finite and in general position, with $\operatorname{mir}(p)\in S$ for every $p\in S$,
and suppose that $S$ has no convex $2j$-gon. Then $\card S\le2M+2$.
\end{theorem}

\begin{proof}
We follow the Lean proof. The set $T=R(S)$ satisfies the assumptions of Lemma~\ref{mm:lem:upper-caps},
with $p_1>0$ for all $p\in T$. Let $C=\mathrm{cfg}(R(S))$. If $C$ had a $j$-cap,
Lemma~\ref{mm:lem:upper-caps}(a) with $m=j\ge2$ would give a convex $2j$-gon in $S$; so $C$ has no
$j$-cap. The set $\mathrm{pts}(C)=R(S)\subseteq S$ has no convex $2j$-gon (Lemma~\ref{fd:lem:mono}(c)).
By (hM), $\card{R(S)}\le M$, and by Lemma~\ref{mm:lem:count},
$\card S=2\card{R(S)}+\card{A(S)}\le2M+2$.
\end{proof}

\begin{remark}[the hypothesis (hM), and indexing]\label{mm:rem:hm}
We comment on the hypothesis of Theorems~\ref{mm:thm:odd} and~\ref{mm:thm:even} and of parts (b), (c) of Theorem~\ref{mm:thm:line}.
\begin{enumerate}[label=(\arabic*)]
\item \emph{(hM) is a parametric hypothesis.} $M$ is any natural number for which (hM) holds; it is not
  defined as a maximum. Theorems~\ref{mm:thm:odd}, \ref{mm:thm:even} and~\ref{mm:thm:line}(b),(c) are
  implications from (hM) to a bound on $\card S$. Thus (hM) is an assumption on the parameter $M$, not an
  unproved claim of the project: for a given pair $(j,M)$ the bound holds as soon as (hM) is proved for that
  pair. In the project, (hM) is discharged by Lean-checked theorems in the following instances. Odd case:
  $(j,M)=(3,16)$ by \lean{MirrorOdd.capacity16} (in \lean{MirrorOdd.mirror_seven'},
  \dirn{MIRRORETV/lean/MirrorMain.lean:2631}) and $(j,M)=(4,113)$ by \lean{C12Check.capacity113} (in
  \lean{C12Check.mirror_nine}, \dirn{MIRRORETV/lean/MirrorNine_append.lean:131}). Even case:
  $(j,M)=(4,22)$ by \lean{MirrorOdd.capacity22g} (in \lean{MirrorOdd.mirror_eight'},
  \dirn{MIRRORETV/lean/MirrorMain.lean:2636}). The axioms of \lean{mirror_seven'} and \lean{mirror_eight'}
  are recorded in \dirn{MIRRORETV/lean/MirrorMain.out} (lines 860--861), those of \lean{C12Check.mirror_nine}
  in \dirn{MIRRORETV/lean/MirrorNine.out} (line 881; that log checks
  \dirn{MIRRORETV/lean/MirrorNine_append.lean} appended to \dirn{MIRRORETV/lean/MirrorMain.lean}); in each
  case they are the three standard axioms. These instances rest on the repository's Lean formalization of
  Baek's theorem (the vendored theorem \lean{main}, Theorem~\ref{mb:thm:baek}), and the instance
  $(4,113)$ also on the vendored cups--caps theorem \lean{Config.cap_cup}. They are treated in the section
  of Theorems~\ref{mb:thm:seven}, \ref{mb:thm:eight} and~\ref{mb:thm:nine}: the three capacity bounds
  appear there as statements used in the proofs of these theorems, \lean{mirror_seven'} and
  \lean{mirror_eight'} in remarks after Theorems~\ref{mb:thm:seven} and~\ref{mb:thm:eight}, and
  \lean{C12Check.mirror_nine} is part~(b) of Theorem~\ref{mb:thm:nine}. For small $j$:
  Lemma~\ref{mm:lem:cup}(b) is exactly (hM) for the odd case
  $(j,M)=(2,4)$ (with $K=5$) and for the even case $(j,M)=(3,5)$ (with $K=6$); and (hM) for the even case
  $(j,M)=(2,1)$ holds because any two points of a configuration form a $2$-cap. So the bounds $8$, $12$
  and $4$ of Theorem~\ref{mm:thm:small} also follow from Theorems~\ref{mm:thm:odd}
  and~\ref{mm:thm:even}; the Lean proofs of Theorem~\ref{mm:thm:small} do not take this route.
\item \emph{Relation to the Erd\H{o}s--Tuza--Valtr function.} \dirn{MIRRORETV/README.md} and
  \dirn{MIRRORETV/PROOFS.md} state Theorem~A with Baek's number $N(n,a,b)$ in the place of $M$: the
  largest size of an $x$-distinct set in general position with no $n$ points in convex position, no
  $a$-cap and no $b$-cup, where caps and cups are defined by consecutive triples. In their own index $j'$
  (see (3)) the odd case uses $N(2j'-1,j',2j'-1)$ and the even case $N(2j',j',2j')$; in the indexing of
  Theorems~\ref{mm:thm:odd} and~\ref{mm:thm:even} these are $N(2j+1,j+1,2j+1)$ (with $j'=j+1$) and
  $N(2j,j,2j)$ (with $j'=j$). These numbers are finite: by the cups--caps theorem of Erd\H{o}s and
  Szekeres (1935), for $a,b\ge2$ every set in the class that defines $N(n,a,b)$ has at most
  $\binom{a+b-4}{a-2}$ points (\dirn{MIRRORETV/PROOFS.md}, \S1), and the two numbers above have
  $a,b\ge j\ge2$.
  That (hM) holds for $M$ equal to the corresponding number then follows, since the point set of a
  configuration of $n$ points is an $x$-distinct set of $n$ points in general position, from two facts
  stated in writing in \dirn{MIRRORETV/PROOFS.md}: consecutive clockwise triples imply that all
  increasing triples are clockwise (Lemma~2.1(b)), and a cup is in convex position (\S1; see also
  Lemma~2.3). No declaration of \lean{MirrorOdd} states this, and it is not used in this section. A
  separate Lean formalization of Theorem~A in this form, written by a review with its own definitions
  and importing only Mathlib, is
  \dirn{MIRRORETV/review/MirrorA.lean} (for example \lean{theoremA_odd_bound}, whose hypothesis is an
  upper bound on the sets of the class that defines $N(2j'-1,j',2j'-1)$). It is not restated in this
  document; the Coverage paragraph at the end of this section describes it and gives the reason.
\item \emph{Index shift.} In Theorem~\ref{mm:thm:odd}, $k=2j+1$ and the caps have $j+1$ points. The
  written sources state the odd case as $k=2j'-1$ with $j'$-caps, so the Lean $j$ equals $j'-1$; and
  $2\le j$ means $k\ge5$, so the case $k=3$ is not covered by Theorem~\ref{mm:thm:odd}. In
  Theorem~\ref{mm:thm:even}, $k=2j$ with $j$-caps, the same $j$ as in the written sources.
\item In the odd case, Step~3 adds an axis point to $R(S)$ at no cost: by
  Lemma~\ref{mm:lem:upper-caps}(b) a $(j+1)$-cap through it would still give a convex $(2j+1)$-gon. This
  yields $2M$ instead of $2M+2$. In the even case a $j$-cap with an endpoint on the axis gives only
  $2j-1$ points, and the (at most two) axis points are simply counted.
\end{enumerate}
\end{remark}

\subsection{Arbitrary mirror lines}

\begin{definition}[line reflection and frame; \leanref{MirrorOdd.lineRefl}{MIRRORETV/lean/MirrorMain.lean:1879}, \leanref{MirrorOdd.frame}{MIRRORETV/lean/MirrorMain.lean:1883}, \leanref{MirrorOdd.unframeLin}{MIRRORETV/lean/MirrorMain.lean:1886}, \leanref{MirrorOdd.unframe}{MIRRORETV/lean/MirrorMain.lean:1892}]\label{mm:def:frame}
Let $n_0,n_1,c\in\R$. For $p,q\in\R^2$ put
\begin{align*}
\mathrm{lineRefl}(p)&=\bigl(p_0-2(n_0p_0+n_1p_1-c)\,n_0,\ \ p_1-2(n_0p_0+n_1p_1-c)\,n_1\bigr),\\
\mathrm{frame}(p)&=\bigl(n_1(p_0-cn_0)-n_0(p_1-cn_1),\ \ n_0(p_0-cn_0)+n_1(p_1-cn_1)\bigr),\\
\mathrm{unframeLin}(q)&=\bigl(n_1q_0+n_0q_1,\ \ -n_0q_0+n_1q_1\bigr),\\
\mathrm{unframe}(q)&=\bigl(n_1q_0+n_0q_1+cn_0,\ \ -n_0q_0+n_1q_1+cn_1\bigr).
\end{align*}
These maps depend on $n_0,n_1,c$ (Lean \lean{lineRefl}, \lean{frame} and \lean{unframe} take the
arguments $n_0,n_1,c$; \lean{unframeLin} takes $n_0,n_1$). The map $\mathrm{unframeLin}$ is linear, and
$\mathrm{unframe}$ is the affine map with linear part $\mathrm{unframeLin}$ and translation part
$(cn_0,cn_1)$.
\end{definition}

\begin{remark}
Let $n_0^2+n_1^2=1$ and $\ell=\{p:n_0p_0+n_1p_1=c\}$; every line in $\R^2$ has this form. With
$\lambda(p)=n_0p_0+n_1p_1-c$ we have $\mathrm{lineRefl}(p)=p-2\lambda(p)(n_0,n_1)$, so
$\lambda(\mathrm{lineRefl}(p))=-\lambda(p)$ while $n_1p_0-n_0p_1$ is unchanged: $\mathrm{lineRefl}$ is the
orthogonal reflection in $\ell$. Moreover $\mathrm{frame}(p)=(n_1p_0-n_0p_1,\ \lambda(p))$ is a rotation by
the matrix $\begin{psmallmatrix}n_1&-n_0\\ n_0&n_1\end{psmallmatrix}$ followed by a translation, and it maps
$\ell$ onto the $x$-axis. This explains the definitions; only Lemma~\ref{mm:lem:frame} is used below.
\end{remark}

\begin{lemma}[frame identities; \leanref{MirrorOdd.unframe_apply}{MIRRORETV/lean/MirrorMain.lean:1897}, \leanref{MirrorOdd.unframe_frame}{MIRRORETV/lean/MirrorMain.lean:1900}, \leanref{MirrorOdd.frame_unframe}{MIRRORETV/lean/MirrorMain.lean:1907}, \leanref{MirrorOdd.frame_lineRefl}{MIRRORETV/lean/MirrorMain.lean:1914}, \leanref{MirrorOdd.orient_frame}{MIRRORETV/lean/MirrorMain.lean:1920}, \leanref{MirrorOdd.frame_injective}{MIRRORETV/lean/MirrorMain.lean:1925}, \leanref{MirrorOdd.unframe_injective}{MIRRORETV/lean/MirrorMain.lean:1930}]\label{mm:lem:frame}
Let $n_0,n_1,c\in\R$. For every $q\in\R^2$, $\mathrm{unframe}(q)$ is given by the formula of
Definition~\ref{mm:def:frame} (\lean{unframe_apply}; no hypothesis). If $n_0^2+n_1^2=1$, then for all
$p,q,a,b,d\in\R^2$:
\begin{enumerate}[label=(\alph*)]
\item $\mathrm{unframe}(\mathrm{frame}(p))=p$ (\lean{unframe_frame});
\item $\mathrm{frame}(\mathrm{unframe}(q))=q$ (\lean{frame_unframe});
\item $\mathrm{frame}(\mathrm{lineRefl}(p))=\operatorname{mir}(\mathrm{frame}(p))$ (\lean{frame_lineRefl});
\item $\orient(\mathrm{frame}(a),\mathrm{frame}(b),\mathrm{frame}(d))=\orient(a,b,d)$ (\lean{orient_frame});
\item $\mathrm{frame}$ and $\mathrm{unframe}$ are injective (\lean{frame_injective},
  \lean{unframe_injective}).
\end{enumerate}
\end{lemma}

\begin{proof}
The formula for $\mathrm{unframe}$ holds by definition (in Lean, \lean{unframe_apply} is proved by
\lean{rfl}). Put $s=n_0^2+n_1^2$ and $\lambda=n_0p_0+n_1p_1-c$. Expanding the
definition gives, with no hypothesis on $n_0,n_1,c$,
\[ \mathrm{frame}(p)=\bigl(n_1p_0-n_0p_1,\ n_0p_0+n_1p_1-cs\bigr)=Rp-(0,cs),\qquad
   R=\begin{pmatrix}n_1&-n_0\\ n_0&n_1\end{pmatrix}. \]
We prove four polynomial identities, valid for all $n_0,n_1,c$.
\begin{itemize}
\item $\mathrm{unframe}(\mathrm{frame}(p))$ has first coordinate
  $n_1(n_1p_0-n_0p_1)+n_0(n_0p_0+n_1p_1-cs)+cn_0=sp_0-csn_0+cn_0=p_0+(s-1)(p_0-cn_0)$, and second coordinate
  $-n_0(n_1p_0-n_0p_1)+n_1(n_0p_0+n_1p_1-cs)+cn_1=sp_1-csn_1+cn_1=p_1+(s-1)(p_1-cn_1)$.
\item With $X=n_1q_0+n_0q_1+cn_0$ and $Y=-n_0q_0+n_1q_1+cn_1$, we get
  $\mathrm{frame}(\mathrm{unframe}(q))=(n_1X-n_0Y,\ n_0X+n_1Y-cs)=(sq_0,\ sq_1+cs-cs)=q+(s-1)q$.
\item $\mathrm{lineRefl}(p)=p-2\lambda\,(n_0,n_1)$ and $R\,(n_0,n_1)^{\mathsf T}=(0,s)^{\mathsf T}$, so
  $\mathrm{frame}(\mathrm{lineRefl}(p))=\mathrm{frame}(p)-2\lambda\,(0,s)$. Its first coordinate equals that
  of $\mathrm{frame}(p)$, which is also that of $\operatorname{mir}(\mathrm{frame}(p))$. Its second
  coordinate minus that of $\operatorname{mir}(\mathrm{frame}(p))$ is
  $2(n_0p_0+n_1p_1-cs)-2\lambda s=-2(s-1)(n_0p_0+n_1p_1)$.
\item $\mathrm{frame}$ is the affine map $x\mapsto Rx-(0,cs)$ with $\det R=s$, so
  $\orient(\mathrm{frame}(a),\mathrm{frame}(b),\mathrm{frame}(d))=s\,\orient(a,b,d)$ by
  Lemma~\ref{fd:lem:affine-orient}.
\end{itemize}
With $s=1$ these identities give (a)--(d). (They are the certificates of the Lean proofs: the first
coordinate in (c) is closed by \lean{ring}, and every other coordinate, and (d), by
\lean{linear_combination} with the multiples $p_0-cn_0$, $p_1-cn_1$, $q_0$, $q_1$, $-2(n_0p_0+n_1p_1)$ and
$\orient(a,b,d)$ of the hypothesis $s=1$.)
(e) By (a), $\mathrm{unframe}$ is a left inverse of $\mathrm{frame}$, so $\mathrm{frame}$ is injective; by
(b), $\mathrm{frame}$ is a left inverse of $\mathrm{unframe}$, so $\mathrm{unframe}$ is injective.
\end{proof}

\begin{theorem}[transport, and Theorem A for an arbitrary mirror line; \leanref{MirrorOdd.transport}{MIRRORETV/lean/MirrorMain.lean:1937}, \leanref{MirrorOdd.mirror_odd_general_line}{MIRRORETV/lean/MirrorMain.lean:2641}, \leanref{MirrorOdd.mirror_even_general_line}{MIRRORETV/lean/MirrorMain.lean:2650}]\label{mm:thm:line}
Part (a) moves a set that is symmetric in a line to a set that is symmetric in the $x$-axis; parts (b)
and~(c) are Theorems~\ref{mm:thm:odd} and~\ref{mm:thm:even} for an arbitrary mirror line.
\begin{enumerate}[label=(\alph*)]
\item (\lean{transport}) Let $n_0,n_1,c\in\R$ with $n_0^2+n_1^2=1$, let $S\subset\R^2$ be finite, and let
  $m\in\N$. Suppose that $S$ is in general position, that $\mathrm{lineRefl}(p)\in S$ for every $p\in S$,
  and that $S$ has no convex $m$-gon. Let $S'=\mathrm{frame}(S)$. Then (i) $S'$ is in general position;
  (ii) $\operatorname{mir}(q)\in S'$ for every $q\in S'$; (iii) $S'$ has no convex $m$-gon; and
  (iv) $\card{S'}=\card S$.
\item (\lean{mirror_odd_general_line}) Let $j,M\in\N$ with $2\le j$, and assume the parametric
  hypothesis (hM) of Theorem~\ref{mm:thm:odd}. Let $n_0,n_1,c\in\R$ with $n_0^2+n_1^2=1$, and let
  $S\subset\R^2$ be finite and in general position, with $\mathrm{lineRefl}(p)\in S$ for every $p\in S$
  and with no convex $(2j+1)$-gon. Then $\card S\le2M$.
\item (\lean{mirror_even_general_line}) Let $j,M\in\N$ with $2\le j$, and assume the parametric
  hypothesis (hM) of Theorem~\ref{mm:thm:even}. Let $n_0,n_1,c\in\R$ with $n_0^2+n_1^2=1$, and let
  $S\subset\R^2$ be finite and in general position, with $\mathrm{lineRefl}(p)\in S$ for every $p\in S$
  and with no convex $2j$-gon. Then $\card S\le2M+2$.
\end{enumerate}
Here $\mathrm{lineRefl}$ and $\mathrm{frame}$ are the maps of Definition~\ref{mm:def:frame} for the given
$n_0,n_1,c$.
\end{theorem}

\begin{proof}
(a) (i) Let $a',b',d'\in S'$ be pairwise distinct, say $a'=\mathrm{frame}(a)$, $b'=\mathrm{frame}(b)$,
$d'=\mathrm{frame}(d)$ with $a,b,d\in S$. Then $a,b,d$ are pairwise distinct (equal points have equal
images), so $\orient(a,b,d)\ne0$ by Lemma~\ref{fd:lem:collinear}, and
$\orient(a',b',d')=\orient(a,b,d)\ne0$ by Lemma~\ref{mm:lem:frame}(d). By Lemma~\ref{fd:lem:collinear},
$S'$ is in general position.
(ii) If $q=\mathrm{frame}(p)$ with $p\in S$, then $\operatorname{mir}(q)=\mathrm{frame}(\mathrm{lineRefl}(p))$
by Lemma~\ref{mm:lem:frame}(c), and $\mathrm{lineRefl}(p)\in S$.
(iii) Suppose that $S'$ has a convex $m$-gon. The map $\mathrm{unframe}$ is affine and injective
(Lemma~\ref{mm:lem:frame}(e)), so $\mathrm{unframe}(S')$ has a convex $m$-gon by
Lemma~\ref{mm:lem:ngon}(b). But $\mathrm{unframe}(S')=\{\mathrm{unframe}(\mathrm{frame}(p)):p\in S\}=S$ by
Lemma~\ref{mm:lem:frame}(a), a contradiction.
(iv) $\mathrm{frame}$ is injective (Lemma~\ref{mm:lem:frame}(e)).

(b) By (a) with $m=2j+1$, the finite set $S'=\mathrm{frame}(S)$ is in general position, is
$\operatorname{mir}$-symmetric, has no convex $(2j+1)$-gon, and $\card{S'}=\card S$.
Theorem~\ref{mm:thm:odd} (with the same $j$, $M$ and (hM)) gives $\card S=\card{S'}\le2M$.

(c) The same with $m=2j$ and Theorem~\ref{mm:thm:even}. This is also the structure of the Lean proofs.
\end{proof}

\subsection{Small \texorpdfstring{$k$}{k}}

\begin{lemma}[configurations without a $3$-cap; \leanref{MirrorOdd.hasCap3_of}{MIRRORETV/lean/MirrorMain.lean:1991}, \leanref{MirrorOdd.capacity_cup}{MIRRORETV/lean/MirrorMain.lean:2002}]\label{mm:lem:cup}
Both parts concern configurations in the sense of Definition~\ref{mm:def:config}.
\begin{enumerate}[label=(\alph*)]
\item (\lean{hasCap3_of}) Let $C$ be a configuration of $n$ points, and let $a<b<d$ in $[n]$ with
  $\orient(C_a,C_b,C_d)<0$. Then $C$ has a $3$-cap.
\item (\lean{capacity_cup}) Let $n,K\in\N$ with $4\le K$, and let $C$ be a configuration of $n$ points
  that has no $3$-cap and such that $\mathrm{pts}(C)$ has no convex $K$-gon. Then $n\le K-1$.
\end{enumerate}
\end{lemma}

\begin{proof}
(a) The map $f:[3]\to[n]$ with $f(0)=a$, $f(1)=b$, $f(2)=d$ is strictly increasing, and the only triple
$i<j<l$ in $[3]$ is $(0,1,2)$, for which the required inequality is the hypothesis.

(b) Suppose $n>K-1$, that is, $n\ge K$. The map $i\mapsto C_i$ is injective because the abscissae
increase strictly. Let $e:[K]\to[n]$ be the inclusion $i\mapsto i$ (Lean \lean{Fin.castLEEmb}). For
$i<j<l$ in $[K]$ we have $\orient(C_i,C_j,C_l)\ne0$ (Definition~\ref{mm:def:config}(a)), and
$\orient(C_i,C_j,C_l)<0$ is impossible because by (a) it would give a $3$-cap; hence
$\orient(C_i,C_j,C_l)>0$. Lemma~\ref{fd:lem:pos-cycle} with $N=n$, $k=K\ge4$, $p=C$, this $e$ and the set
$\mathrm{pts}(C)$ (which contains every $C_{e(i)}$) gives a convex $K$-gon in $\mathrm{pts}(C)$, a
contradiction.
\end{proof}

\begin{theorem}[$k=4,5,6$; \leanref{MirrorOdd.mirror_four}{MIRRORETV/lean/MirrorMain.lean:2019}, \leanref{MirrorOdd.mirror_five}{MIRRORETV/lean/MirrorMain.lean:2040}, \leanref{MirrorOdd.mirror_six}{MIRRORETV/lean/MirrorMain.lean:2129}, \leanref{MirrorOdd.mirror_four_line}{MIRRORETV/lean/MirrorMain.lean:2159}, \leanref{MirrorOdd.mirror_five_line}{MIRRORETV/lean/MirrorMain.lean:2165}, \leanref{MirrorOdd.mirror_six_line}{MIRRORETV/lean/MirrorMain.lean:2171}]\label{mm:thm:small}
In (a)--(c) let $S\subset\R^2$ be finite and in general position, with $\operatorname{mir}(p)\in S$ for
every $p\in S$. In (d)--(f) let $n_0,n_1,c\in\R$ with $n_0^2+n_1^2=1$, and let $S\subset\R^2$ be finite
and in general position, with $\mathrm{lineRefl}(p)\in S$ for every $p\in S$.
\begin{enumerate}[label=(\alph*)]
\item (\lean{mirror_four}) If $S$ has no convex $4$-gon, then $\card S\le4$.
\item (\lean{mirror_five}) If $S$ has no convex $5$-gon, then $\card S\le8$.
\item (\lean{mirror_six}) If $S$ has no convex $6$-gon, then $\card S\le12$.
\item (\lean{mirror_four_line}) If $S$ has no convex $4$-gon, then $\card S\le4$.
\item (\lean{mirror_five_line}) If $S$ has no convex $5$-gon, then $\card S\le8$.
\item (\lean{mirror_six_line}) If $S$ has no convex $6$-gon, then $\card S\le12$.
\end{enumerate}
These statements have no hypotheses beyond the displayed ones.
\end{theorem}

\begin{proof}
We follow the Lean proofs.

(a) First, $\card{R(S)}\le1$. Otherwise there are $p\ne q$ in $R(S)$. Lemma~\ref{mm:lem:injx}, applied with
$T=R(S)$ (all its points have positive ordinate, so the condition on axis points is vacuous), gives
$p_0\ne q_0$; exchanging the names if necessary, $p_0<q_0$. Proposition~\ref{mm:prop:small-polygons}(a)
with $(c_0,c_1)=(p,q)$ and $P=S$ applies ($p_1,q_1>0$, and $\operatorname{mir}p,\operatorname{mir}q,q,p\in S$)
and gives a convex $4$-gon in $S$, a contradiction. By Lemma~\ref{mm:lem:count},
$\card S=2\card{R(S)}+\card{A(S)}\le2+2=4$.

(b) \emph{Step~1:} if $T\subseteq S$ has $p_1\ge0$ for all $p\in T$ and at most one point on the axis,
then $\card T\le4$. Let $C=\mathrm{cfg}(T)$ (Lemma~\ref{mm:lem:upper-caps}). $C$ has no $3$-cap, because
otherwise Lemma~\ref{mm:lem:upper-caps}(b) with $m=3$ gives a convex $5$-gon in $S$. (The Lean proof makes
the same case distinction as the proof of Lemma~\ref{mm:lem:upper-caps}(b) for the cap $c_0,c_1,c_2$, but
applies Proposition~\ref{mm:prop:small-polygons}(b) when the three points lie strictly above the axis;
Proposition~\ref{mm:prop:small-polygons}(g), with the points $a,c_0,c_1$ of the Proposition taken to be
$c_2,c_0,c_1$, when $c_2$ lies on the axis; and Proposition~\ref{mm:prop:small-polygons}(f), with
$a,c_0,c_1$ taken to be $c_0,c_1,c_2$, when $c_0$ lies on the axis. The case that $c_1$ lies on the axis
is excluded by Lemma~\ref{mm:lem:slope}(b).) Moreover $\mathrm{pts}(C)=T\subseteq S$ has no convex $5$-gon.
Lemma~\ref{mm:lem:cup}(b) with $K=5$ gives $\card T\le4$.
\emph{Step~2:} $\card{R(S)}\le4$ by Step~1.
\emph{Step~3:} if $A(S)=\emptyset$, then $\card S=2\card{R(S)}\le8$. Otherwise, for $a\in A(S)$, Step~1
applied to $R(S)\cup\{a\}$ (as in Step~3 of Theorem~\ref{mm:thm:odd}) gives $\card{R(S)}+1\le4$, and
$\card S\le2\cdot3+2=8$ by Lemma~\ref{mm:lem:count}.

(c) Let $C=\mathrm{cfg}(R(S))$ (Lemma~\ref{mm:lem:upper-caps} with $T=R(S)$). Suppose $C$ has a $3$-cap,
given by $f$, and put $c_i=C_{f(i)}$; then $(c_0)_0<(c_1)_0<(c_2)_0$, all $(c_i)_1>0$,
$\orient(c_0,c_1,c_2)<0$, and $S$ contains $c_i$ and $\operatorname{mir}c_i$. Then
Proposition~\ref{mm:prop:small-polygons}(c) with $P=S$ (equivalently, Lemma~\ref{mm:lem:upper-caps}(a)
with $m=3$) gives a convex $6$-gon in $S$, a contradiction. Also $\mathrm{pts}(C)=R(S)\subseteq S$ has no
convex $6$-gon. Lemma~\ref{mm:lem:cup}(b) with $K=6$ gives $\card{R(S)}\le5$, so
$\card S\le2\cdot5+2=12$ by Lemma~\ref{mm:lem:count}.

(d)--(f) By Theorem~\ref{mm:thm:line}(a) with $m=4,5,6$, the set $S'=\mathrm{frame}(S)$ satisfies the
hypotheses of (a)--(c), respectively, and $\card{S'}=\card S$.
\end{proof}

\paragraph{Coverage.}
{\raggedright
Covered, all in namespace \lean{MirrorOdd} of \dirn{MIRRORETV/lean/MirrorMain.lean} unless stated
otherwise:
\lean{mir}, \lean{Rset}, \lean{Aset} (Definition~\ref{mm:def:mir});
\lean{mir_zero}, \lean{mir_one}, \lean{mir_mir}, \lean{orient_mir3}, \lean{orient_mir2},
\lean{orient_cycl}, \lean{orient_swap12}, \lean{orient_rev} (Lemma~\ref{mm:lem:orient-mir});
\lean{ErdosSzekeres.Omission.Configuration} and \lean{Configuration.HasCap}
(\dirn{ErdosSzekeres/Omission/Realization.lean}), \lean{xset}, \lean{emb}, \lean{sortPt}, \lean{cfg}
(Definition~\ref{mm:def:config});
\lean{xset_card}, \lean{emb_mem}, \lean{sortPt_mem}, \lean{sortPt_x}, \lean{sortPt_strictMono},
\lean{sortPt_injective}, \lean{sortPt_image} (Lemma~\ref{mm:lem:sort});
\lean{hasConvexNGon_of_le}, \lean{hasConvexNGon_affineMap} (Lemma~\ref{mm:lem:ngon});
\lean{cap_below_lines} (Lemma~\ref{mm:lem:slope});
\lean{cap_mirror_polygon} (Theorem~\ref{mm:thm:mirror-polygon});
\lean{axis_left_polygon}, \lean{axis_right_polygon} (Theorem~\ref{mm:thm:axis-polygon});
\lean{cyc4}--\lean{cyc8} (Lemma~\ref{mm:lem:cyc});
\lean{cap_mirror_quadrilateral}, \lean{cap_mirror_pentagon}, \lean{cap_mirror_hexagon},
\lean{cap_mirror_heptagon}, \lean{cap_mirror_octagon}, \lean{axis_left_pentagon},
\lean{axis_right_pentagon}, \lean{axis_left_heptagon}, \lean{axis_right_heptagon}
(Proposition~\ref{mm:prop:small-polygons});
\lean{card_split}, \lean{card_A_le} (Lemma~\ref{mm:lem:count}); \lean{injOn_x}
(Lemma~\ref{mm:lem:injx});
\lean{mirror_odd_general} (Theorem~\ref{mm:thm:odd}); \lean{mirror_even_general}
(Theorem~\ref{mm:thm:even});
\lean{lineRefl}, \lean{frame}, \lean{unframeLin}, \lean{unframe} (Definition~\ref{mm:def:frame});
\lean{unframe_apply}, \lean{unframe_frame}, \lean{frame_unframe}, \lean{frame_lineRefl},
\lean{orient_frame}, \lean{frame_injective}, \lean{unframe_injective} (Lemma~\ref{mm:lem:frame});
\lean{transport}, \lean{mirror_odd_general_line}, \lean{mirror_even_general_line} (Theorem~\ref{mm:thm:line});
\lean{hasCap3_of}, \lean{capacity_cup} (Lemma~\ref{mm:lem:cup});
\lean{mirror_four}, \lean{mirror_five}, \lean{mirror_six} and their \lean{_line} versions
(Theorem~\ref{mm:thm:small}).
Three items are written steps with no Lean declaration of their own: Lemma~\ref{mm:lem:upper-caps}
isolates the cap argument of the local claims \lean{keyT} and \lean{hR}; part~(b) of
Lemma~\ref{mm:lem:slope} states the local claim \lean{interior} and the inline interior-point argument of
the Lean proofs; and part~(e) of Lemma~\ref{mm:lem:orient-mir} is an identity that the Lean proof of
\lean{cap_below_lines} uses implicitly when it expands $\orient$ and closes each case by
\lean{nlinarith}.

Not covered here, deliberately, from \dirn{MIRRORETV/lean/MirrorMain.lean}:
\begin{itemize}
\item \lean{hasNGon_heptagon} (line 1512), \lean{capacity16} (1529), \lean{mirror_seven_of} (1676),
  \lean{mirror_seven} (1803), \lean{capacity22g} (1812), \lean{mirror_eight} (1820),
  \lean{mirror_seven_line} (1969) and \lean{mirror_eight_line} (1977). \lean{hasNGon_heptagon} is the
  bridge, used by \lean{capacity16}, from a $7$-gon of the combinatorial model of a configuration
  (\lean{colorConfig}) to a convex $7$-gon of its point set. The others rest on Baek's theorem (the
  vendored theorem \lean{main}): \lean{capacity16} applies it directly, \lean{capacity22g} through
  \lean{geometric_capacity22}, and the remaining five through these two. They are treated in the section
  of Theorems~\ref{mb:thm:seven} and~\ref{mb:thm:eight}: \lean{mirror_seven}, \lean{mirror_eight} and
  their \lean{_line} versions are these two theorems, and the other four are statements used in their
  proofs.
\item \lean{mirror_seven'} (2631) and \lean{mirror_eight'} (2636): consistency re-derivations of the bounds
  $32$ and $46$. They are Theorem~\ref{mm:thm:odd} with $(j,M)=(3,16)$ and Theorem~\ref{mm:thm:even} with
  $(j,M)=(4,22)$, with (hM) supplied by \lean{capacity16} and \lean{capacity22g} (see
  Remark~\ref{mm:rem:hm}); the bounds themselves are Theorems~\ref{mb:thm:seven}
  and~\ref{mb:thm:eight}.
\item The \texttt{\#print axioms} lines 2662--2677: audit commands; their output is the axioms evidence
  cited at the beginning of this section.
\end{itemize}
Not covered here, deliberately, from the other Lean files of \dirn{MIRRORETV}:
\begin{itemize}
\item The five declarations of \dirn{MIRRORETV/lean/MirrorNine_append.lean} (namespace
  \lean{C12Check}). They are treated in the section of Theorem~\ref{mb:thm:nine}, and
  \lean{config_capacity113}, \lean{mirror_nine} and \lean{mirror_nine_line} are that theorem.
\item The independent formalizations \dirn{MIRRORETV/review/MirrorA.lean} and
  \dirn{MIRRORETV/review/MirrorCap.lean}, written by reviews. Each imports only Mathlib and has its own
  definitions, and no Lean file of the project imports either of them. No statement or proof of this
  document uses them. Their content overlaps that of
  Theorems~\ref{mm:thm:mirror-polygon}, \ref{mm:thm:axis-polygon}, \ref{mm:thm:odd}, \ref{mm:thm:even}
  and~\ref{mm:thm:line}, but in another formulation: points in $\R\times\R$, caps and cups defined by
  consecutive triples, and an upper bound on the class that defines $N(n,a,b)$
  (Remark~\ref{mm:rem:hm}(2)) in place of (hM). Beyond this overlap they contain the odd case $k=3$,
  which Theorem~\ref{mm:thm:odd} excludes (Remark~\ref{mm:rem:hm}(3)); the mirror lemma for caps with
  both endpoints on the axis; the convex position of caps and of cups; and the coincidence statement
  \lean{MirrorCap.coincidences}.
  \begin{itemize}
  \item \dirn{MirrorA.lean} (891 lines, \texttt{1a8a62cb4385}\dots; namespace \lean{MirrorA}; 16
    definitions, 55 theorems and lemmas, and one \texttt{example}). The definitions are \lean{orient}, the
    reflection \lean{mu} in the $x$-axis and its linear form \lean{muL}, \lean{ConvexPos}, the caps and cups
    \lean{IsCapSeq} and \lean{IsCupSeq}, the set \lean{V} of the points of a cap and their mirror images,
    \lean{GenPos}, \lean{XDistinct}, \lean{HasGon}, \lean{HasCap}, \lean{HasCup}, the class \lean{InClass}
    that defines $N(n,a,b)$, \lean{MirrorSym}, \lean{refl}, the reflection in the line through a point
    with a given unit direction, and \lean{isoT}, the rigid motion that maps this line onto the $x$-axis.
    The main results are \lean{lemma31} (282: a cap, in the
    sense of \lean{IsCapSeq}, all of whose points have ordinates $\ge0$, together with its mirror image,
    is in convex position; the review's form of \dirn{MIRRORETV/PROOFS.md}, Lemma~3.1),
    \lean{cap_convexPos} (352) and \lean{cup_convexPos} (417) (the points of a cap, resp.\ of a cup, are in
    convex position), \lean{theoremA_even} (593) and \lean{theoremA_odd} (615) (in the indexing of the
    written sources, for $j\ge2$: if $S$ is symmetric in the $x$-axis, in general position, and has no $2j$, resp.\
    $2j-1$, points in convex position, then some set $X$ in the class of $N(2j,j,2j)$, resp.\
    $N(2j-1,j,2j-1)$, satisfies $\card S\le2\card X+2$, resp.\ $\card S\le2\card X$),
    \lean{theoremA_even_bound} (685) and \lean{theoremA_odd_bound} (694) (the bounds $2N_b+2$ and $2N_b$
    from an upper bound $N_b$ on that class), \lean{transport} (771; the transport from an arbitrary
    mirror line to the $x$-axis by the rigid motion \lean{isoT}), and \lean{theoremA_even_line} (810)
    and \lean{theoremA_odd_line} (821) (the two bounds for the reflection in an arbitrary line). With
    $j=2$ the odd case is $k=3$. The other 45 lemmas are auxiliary, among them three checks of the
    definition of \lean{refl} (\lean{refl_refl}, \lean{refl_fix_line}, \lean{refl_normal}); the
    \texttt{example} (852) checks that \lean{ConvexPos} is not vacuous. The recorded log
    \dirn{MIRRORETV/review/MirrorA.out} prints the three standard axioms for eight of these theorems.
  \item \dirn{MirrorCap.lean} (354 lines, \texttt{a5c6dadb5ac6}\dots; namespace \lean{MirrorCap}; 16
    declarations): the definitions \lean{orient}, \lean{sl} (the slope of an edge) and \lean{V} (the
    points of the cap and their mirror images), ten auxiliary lemmas and theorems (\lean{edge_eq},
    \lean{orient_eq_slopes}, \lean{slope_dec_of_orient}, \lean{sl_anti}, \lean{x_mono}, \lean{chord_right},
    \lean{chord_left}, \lean{y_nonneg}, \lean{exposed_c}, \lean{exposed_mu}), and \lean{convex_position}
    (292), \lean{convexIndependent} (316) and \lean{coincidences} (327). For a cap $c_0,\dots,c_{m-1}$ given
    by coordinates (strictly increasing abscissae, consecutive triples clockwise) whose interior points
    have positive ordinates and whose endpoints have nonnegative ordinates (positive if $m=1$),
    \lean{convex_position} says that no point of $V$ lies in the convex hull of the other points of $V$
    (the form of \dirn{MIRRORETV/PROOFS.md}, Lemma~3.1), and \lean{convexIndependent} says the same in
    Mathlib's terms. From the monotonicity of the abscissae and the positivity of the interior ordinates,
    \lean{coincidences} derives that $c_i=c_k$ or $\operatorname{mir}c_i=\operatorname{mir}c_k$ forces
    $i=k$, and that $c_i=\operatorname{mir}c_k$ forces $i=k$, $(c_i)_1=0$ and $i\in\{0,m-1\}$; by the
    file's comment, this gives $\card V=2m-e$, where $e$ is the number of endpoints of the cap on the axis.
    The recorded log \dirn{MIRRORETV/review/MirrorCap.out} prints the three standard axioms for these three
    theorems.
  \end{itemize}
\item The far-flat files \dirn{MIRRORETV/lean/FarFlat.lean} (\texttt{08a79f881eea}\dots) and
  \dirn{MIRRORETV/review/FarFlatExtra.lean} (144 lines, \texttt{d30a9b874ba5}\dots). Both import only
  Mathlib, and no Lean file of the project imports them. They hold estimates for the far-flat lower-bound
  construction (Theorem~B of \dirn{MIRRORETV/README.md}). This document treats no part of Theorem~B, so
  no section restates either file.
  \begin{itemize}
  \item \dirn{FarFlat.lean}: the root-namespace definition \lean{orient} on six reals (line 12),
    \lean{orient_wzq_neg} (15), \lean{orient_uqz_neg} (34), \lean{orient_cyclic} (55), \lean{orient_sum}
    (60), \lean{inside_triangle} (67) and \lean{line_misses} (89): three orientation estimates
    (\lean{orient_wzq_neg}, \lean{orient_uqz_neg}, \lean{line_misses}), two identities
    (\lean{orient_cyclic}, \lean{orient_sum}) and a barycentric representation (\lean{inside_triangle}).
    The recorded log \dirn{MIRRORETV/lean/FarFlat.out} prints the axioms of five of these six theorems
    (all but \lean{orient_cyclic}), and they are the standard three.
  \item \dirn{FarFlatExtra.lean} (namespace \lean{FFX}): the definition \lean{orient} on six reals (line
    20), \lean{orient_sum} (22), \lean{bary_of_neg} (28), \lean{bary_of_pos} (44), \lean{gp_cross_axis}
    (60), \lean{a1_inside} (69) and \lean{lemma51_core} (91). Here \lean{orient_sum} and
    \lean{bary_of_neg} have the same statements as \lean{orient_sum} and \lean{inside_triangle} of
    \dirn{FarFlat.lean}; \lean{bary_of_pos} is the barycentric representation for three positive
    orientations; \lean{gp_cross_axis} is a general-position estimate for a point above the axis, a point
    below it and a point of the axis to the right of both; \lean{a1_inside} is the counting step of
    Theorem~B; and \lean{lemma51_core} is the composed core of Lemma~5.1 of \dirn{MIRRORETV/PROOFS.md}.
    The recorded log \dirn{MIRRORETV/review/FarFlatExtra.out} prints the axioms of five of these six
    theorems (all but \lean{orient_sum}), the standard three, and four warnings about unused variables.
  \end{itemize}
\item \dirn{MIRRORETV/review/ReplayRepo.lean} (118 lines, \texttt{20e479e8b3ab}\dots) contains no
  definition or theorem, only a command \texttt{\#replay\_repo}, defined by \texttt{elab} at line~51 and
  run at line~118. By its header comment, it imports only library modules, reads the compiled repository
  modules without importing them, and sends the repository constants in the dependency closure of those
  that \dirn{MIRRORETV/lean/MirrorMain.lean} references directly to Lean's kernel again. The recorded log
  \dirn{MIRRORETV/review/ReplayRepo.out} reports $497$ such constants, none of them an axiom, and a
  successful replay. It is an audit and states no theorem, so there is nothing to restate.
\end{itemize}
\par}

% End of the local typesetting group opened after the section heading.
\endgroup
